\documentclass[11pt]{article}
\usepackage[margin=1in]{geometry}
\usepackage[T1]{fontenc}
\usepackage[utf8]{inputenc}
\usepackage[french]{babel}
\usepackage{amsmath,amssymb,amsthm}
\usepackage[colorlinks=true,linkcolor=blue,urlcolor=blue]{hyperref}
\newtheorem{problem}{Problème}
\newenvironment{refs}{\par\small\noindent\emph{Références :}\begin{itemize}\setlength\itemsep{0pt}}{\end{itemize}\normalsize}
\title{Hors de la Caverne \\ \Large Logique, théorie des ensembles et fondements}
\author{}
\date{}
\begin{document}
\maketitle
\noindent\emph{Des problèmes, pas de réponses.} Chaque problème ci-dessous est posé
sans solution. Les références indiquent où trouver les connaissances nécessaires.
\medskip


\section*{Lycée}

\begin{problem}[{Chevaliers et valets --- Lycée}]
\textbf{FND-01.} \textbf{Problème.} Sur une certaine île, tout habitant est un \textit{chevalier},
qui dit toujours la vérité, ou un \textit{valet}, qui ment toujours.
\begin{enumerate}
\item[(a)] Vous rencontrez A et B. A déclare :
\og{} Au moins l'un de nous deux est un valet. \fg{} Déterminez, avec démonstration, ce que sont A et B.
\item[(b)] Vous rencontrez C et D.
C déclare : \og{} Nous sommes tous les deux des valets. \fg{} Déterminez ce que sont C et D.
\item[(c)] Vous arrivez à une bifurcation ; une
branche mène au village, l'autre au marécage. Un habitant isolé, de type inconnu, se tient
à la bifurcation. Trouvez une unique question fermée (oui/non) que vous puissiez poser de sorte que
la réponse vous indique la bonne route, et démontrez que votre question fonctionne, que
l'habitant soit chevalier ou valet.
\end{enumerate}

\end{problem}

\noindent Raymond Smullyan a rendu ces énigmes célèbres --- et forgé les noms de chevalier
et de valet --- dans son recueil de 1978 \textit{What Is the Name of This Book?} C'est de la logique
propositionnelle déguisée : chaque énigme revient à décider quelles attributions de valeurs de
vérité satisfont une petite formule, et (c) est un premier avant-goût de la conception d'une
formule dont la valeur de vérité est invariante sous le comportement d'un adversaire. Les
résoudre par une analyse de cas systématique, plutôt que par un éclair d'intuition, est le
premier pas vers la démonstration.

\begin{refs}
\item Contexte : Chevaliers et valets --- Wikipédia (en anglais) (\url{https://en.wikipedia.org/wiki/Knights_and_Knaves})
\item Lecture : Raymond Smullyan, \textit{What Is the Name of This Book?}
\end{refs}

\bigskip

\begin{problem}[{Seize fonctions de vérité --- Lycée}]
\textbf{FND-02.} \textbf{Problème.}
\begin{enumerate}
\item[(a)] Écrivez les tables de vérité de $ \neg P $,
$ P \wedge Q $, $ P \vee Q $, $ P \to Q $ et $ P \leftrightarrow Q $, et vérifiez par
table de vérité que $ ((P \to Q) \wedge (Q \to R)) \to (P \to R) $ est une tautologie (vraie
sous toute attribution de valeurs).
\item[(b)] Montrez que $ P \to Q $ n'est \textit{pas} logiquement équivalente à sa
réciproque $ Q \to P $, mais qu'elle l'\textit{est} à sa contraposée
$ \neg Q \to \neg P $.
\item[(c)] Comptez combien il existe de fonctions de vérité distinctes de deux variables,
et démontrez que chacune d'elles peut s'exprimer à l'aide des seuls connecteurs $ \neg $
et $ \wedge $.
\end{enumerate}

\end{problem}

\noindent Les tables de vérité ont pris leur forme reconnaissable chez Emil Post et
Ludwig Wittgenstein vers 1920, avec des racines chez Frege, Peirce et Boole. La partie (c) est un
premier théorème de \textit{complétude fonctionnelle} --- le constat qu'un tout petit jeu de
connecteurs engendre toute la logique propositionnelle. Henry Sheffer a montré en 1913 qu'un seul
connecteur (la barre de Sheffer, NON-ET) suffit, fait câblé en dur dans chaque puce
d'ordinateur.

\begin{refs}
\item Contexte : Table de vérité --- Wikipédia (\url{https://fr.wikipedia.org/wiki/Table_de_v%C3%A9rit%C3%A9})
\item Contexte : Complétude fonctionnelle --- Wikipédia (en anglais) (\url{https://en.wikipedia.org/wiki/Functional_completeness})
\item Lecture : Daniel Velleman, \textit{How to Prove It: A Structured Approach}, ch. 1
\end{refs}

\bigskip

\begin{problem}[{Le jeu de la négation des quantificateurs --- Lycée}]
\textbf{FND-03.} \textbf{Problème.} Toutes les variables parcourent les entiers relatifs $\mathbb{Z}$.
\begin{enumerate}
\item[(a)] Écrivez la
négation de chacun des énoncés ci-dessous sous une forme où aucun signe de négation ne précède
un quantificateur, puis déterminez --- avec démonstration --- lesquels des six énoncés (les trois
originaux, les trois négations) sont vrais : (i) $ \forall x\, \exists y\, (y \cdot y = x) $ ;
(ii) $ \exists y\, \forall x\, (x + y = x) $ ; (iii) tout nombre pair supérieur à 2 est
la somme de deux nombres premiers.
\item[(b)] Exhibez une propriété $ P(x,y) $ des couples d'entiers pour
laquelle $ \forall x\, \exists y\, P(x,y) $ est vraie mais $ \exists y\, \forall x\, P(x,y) $
est fausse, et démontrez que, pour toute propriété $ P $, le second énoncé implique le premier.
\end{enumerate}

\end{problem}

\noindent Apprendre que \og{} pour tout x il existe un y \fg{} et \og{} il existe un y pour tout x \fg{}
ne disent pas la même chose est le pas le plus important du calcul vers l'analyse : les
définitions en $\varepsilon$--$\delta$ de Weierstrass vivent tout entières dans les alternances de quantificateurs.
L'énoncé (iii) est la conjecture de Goldbach (1742), toujours non démontrée --- et c'est
exactement pourquoi on vous demande seulement de la nier correctement, non de la trancher ;
voyez la page des Problèmes ouverts. Depuis Hintikka, les logiciens voient les quantificateurs
alternés comme un jeu entre un vérificateur et un falsificateur.

\begin{refs}
\item Contexte : Quantification (logique) --- Wikipédia (\url{https://fr.wikipedia.org/wiki/Quantification_(logique)})
\item Contexte : Calcul des prédicats (logique du premier ordre) --- Wikipédia (\url{https://fr.wikipedia.org/wiki/Calcul_des_pr%C3%A9dicats})
\item Lecture : Daniel Velleman, \textit{How to Prove It: A Structured Approach}, ch. 2
\end{refs}

\bigskip

\begin{problem}[{La récurrence, et une démonstration qui démontre trop --- Lycée}]
\textbf{FND-04.} \textbf{Problème.}
\begin{enumerate}
\item[(a)] Démontrez par récurrence que
$ 1 + 2 + \cdots + n = \frac{n(n+1)}{2} $ et que
$ 1^3 + 2^3 + \cdots + n^3 = (1 + 2 + \cdots + n)^2 $.
\item[(b)] Démontrez par récurrence forte que
tout entier $ n \ge 2 $ est un produit de nombres premiers.
\item[(c)] Voici une célèbre \og{} démonstration \fg{} que tous
les chevaux ont la même couleur. \textit{Cas de base :} tout ensemble de 1 cheval n'a qu'une couleur.
\textit{Pas de récurrence :} étant donné $ n+1 $ chevaux, retirez-en un ; les $ n $ restants
partagent une couleur par hypothèse. Remettez-le, puis retirez-en un autre ; de nouveau les $ n $
restants partagent une couleur. Les deux ensembles se chevauchant, les $ n+1 $ chevaux doivent
tous partager une même couleur. Identifiez le point exact où cet argument échoue, et expliquez
pourquoi.
\end{enumerate}

\end{problem}

\noindent La récurrence a été utilisée implicitement pendant des siècles, énoncée
clairement par Pascal (1654), nommée par De Morgan (1838), et érigée en axiome de l'arithmétique
par Peano (1889). Le sophisme des chevaux, popularisé par George Pólya comme mise en garde,
enseigne la leçon que tout mathématicien en exercice finit par apprendre à ses dépens : une
récurrence ne vaut que ce que vaut le cas où le pas de récurrence s'enclenche pour la première
fois.

\begin{refs}
\item Contexte : Raisonnement par récurrence --- Wikipédia (\url{https://fr.wikipedia.org/wiki/Raisonnement_par_r%C3%A9currence})
\item Contexte : Axiomes de Peano --- Wikipédia (\url{https://fr.wikipedia.org/wiki/Axiomes_de_Peano})
\item Lecture : Daniel Velleman, \textit{How to Prove It: A Structured Approach}, ch. 6
\end{refs}

\bigskip

\begin{problem}[{Le barbier, et l'ensemble de tous les ensembles --- Lycée}]
\textbf{FND-05.} \textbf{Problème.}
\begin{enumerate}
\item[(a)] Dans un certain village, le barbier rase exactement les
villageois qui ne se rasent pas eux-mêmes. Démontrez qu'un tel barbier ne peut pas exister.
\item[(b)] La théorie naïve des ensembles
autorise la formation de l'ensemble $ \{x : P(x)\} $ pour toute propriété $ P $. Formez
$ R = \{x : x \notin x\} $ et déduisez-en une contradiction.
\item[(c)] Expliquez précisément quelle étape
de formation de (b) les axiomes de Zermelo--Fraenkel interdisent, énoncez le schéma d'axiomes de
séparation (compréhension) qui la remplace, et servez-vous de la séparation ainsi que de votre
argument de (b) pour démontrer qu'il n'existe pas d'ensemble de tous les ensembles.
\end{enumerate}

\end{problem}

\noindent Bertrand Russell a trouvé ce paradoxe en 1901 et l'a communiqué à Frege en
1902, au moment précis où le second volume de l'œuvre d'une vie de Frege, les
\textit{Grundgesetze}, était sous presse ; l'appendice où Frege reconnaît l'effondrement est l'un
des documents les plus poignants de la mathématique. Zermelo avait remarqué le paradoxe de son
côté. La réparation --- les axiomes de Zermelo de 1908, étendus par Fraenkel --- est la théorie des
ensembles (ZFC) dans laquelle on fait des mathématiques aujourd'hui.

\begin{refs}
\item Contexte : Paradoxe de Russell --- Wikipédia (\url{https://fr.wikipedia.org/wiki/Paradoxe_de_Russell})
\item Contexte : Paradoxe du barbier --- Wikipédia (\url{https://fr.wikipedia.org/wiki/Paradoxe_du_barbier})
\item Contexte : Théorie des ensembles de Zermelo--Fraenkel --- Wikipédia (\url{https://fr.wikipedia.org/wiki/Th%C3%A9orie_des_ensembles_de_Zermelo-Fraenkel})
\item Lecture : Paul Halmos, \textit{Naive Set Theory}
\end{refs}

\bigskip

\begin{problem}[{Une phrase qu'aucun insulaire ne peut prononcer --- Lycée, short}]
\textbf{FND-24.} \textbf{Problème.} Sur l'île des chevaliers et des valets de FND-01, démontrez
qu'aucun habitant --- chevalier ou valet --- ne peut jamais dire \og{} Je suis un valet \fg{}.
\end{problem}

\noindent C'est le paradoxe du menteur désamorcé : la phrase \og{} Je mens \fg{}, attribuée dans
l'Antiquité à Épiménide le Crétois, ne peut être affirmée de façon cohérente par quiconque est
astreint à la vérité ou au mensonge. Smullyan a fait de cette observation même le moteur de ses
recueils d'énigmes --- quand un personnage dit quelque chose qu'aucun insulaire ne pourrait dire,
ce sont les hypothèses elles-mêmes qui doivent céder.

\begin{refs}
\item Contexte : Paradoxe du menteur --- Wikipédia (\url{https://fr.wikipedia.org/wiki/Paradoxe_du_menteur})
\item Lecture : Raymond Smullyan, \textit{What Is the Name of This Book?}
\end{refs}

\bigskip

\begin{problem}[{Compter les parties --- Lycée, short}]
\textbf{FND-25.} \textbf{Problème.} Démontrez qu'un ensemble à $ n $ éléments possède
exactement $ 2^n $ parties, en comptant l'ensemble vide et l'ensemble tout entier.
\end{problem}

\noindent La collection de toutes les parties de $ A $ est l'\textit{ensemble des parties}
$ \mathcal{P}(A) $, et ce dénombrement est la raison de la notation $ 2^A $ : une partie,
c'est la même chose qu'une réponse par oui ou non pour chaque élément. Le théorème de Cantor
(FND-07), c'est ce qu'il advient de ce fait innocent quand $ n $ devient infini --- l'ensemble
des parties dépasse l'ensemble, à jamais.

\begin{refs}
\item Contexte : Ensemble des parties d'un ensemble --- Wikipédia (\url{https://fr.wikipedia.org/wiki/Ensemble_des_parties_d%27un_ensemble})
\item Lecture : Daniel Velleman, \textit{How to Prove It: A Structured Approach}, ch. 6
\end{refs}

\bigskip

\begin{problem}[{L'hôtel de Hilbert --- Lycée}]
\textbf{FND-36.} \textbf{Problème.} Un hôtel a des chambres numérotées $ 1, 2, 3, \ldots $ --- une par entier
strictement positif --- et toutes les chambres sont occupées. Dans chaque partie ci-dessous, donnez
une règle explicite de réattribution des clients aux chambres et \textit{démontrez} que votre règle
est une bijection, de sorte que personne ne soit expulsé et que personne ne partage sa chambre.
\begin{enumerate}
\item[(a)] Un nouveau client arrive. Logez tout le monde.
\item[(b)] Un car transportant une infinité de nouveaux clients, numérotés $ 1, 2, 3, \ldots $,
arrive. Logez tout le monde.
\item[(c)] Une infinité de cars arrivent, chacun transportant une infinité de clients, ceux-ci étant
repérés par des couples (numéro du car, numéro du siège). Logez tout le monde.
\item[(d)] On dit qu'un ensemble $ A $ est \textit{infini au sens de Dedekind} s'il existe une bijection
entre $ A $ et une partie propre de $ A $. Démontrez que $ \mathbb{N} $ est infini au sens
de Dedekind, qu'aucun ensemble fini ne l'est, et que tout ensemble infini au sens de Dedekind
contient une partie en bijection avec $ \mathbb{N} $.
\end{enumerate}

\end{problem}

\noindent Hilbert s'est servi de cet hôtel dans une conférence de 1924--25, \og{} Über das
Unendliche \fg{}, pour faire voir que \og{} l'hôtel est plein \fg{} et \og{} on ne peut plus loger personne \fg{}
cessent d'être la même phrase dès que les chambres sont en nombre infini ; George Gamow l'a porté
devant le grand public dans \textit{One Two Three\dots{} Infinity} (1947). La mathématique est plus
ancienne et plus tranchante que la plaisanterie : Dedekind avait déjà pris la propriété de
l'hôtel comme \textit{définition} de l'infini dans \textit{Was sind und was sollen die Zahlen?}
(1888) --- un ensemble est infini quand on peut le mettre en correspondance avec une partie de
lui-même. La partie (d) cache une subtilité qu'il vaut la peine de connaître : démontrer que tout
ensemble infini est infini au sens de Dedekind exige une forme faible de l'axiome du choix
(FND-09), et dans ZF seul les deux notions se séparent.

\begin{refs}
\item Contexte : Hôtel de Hilbert --- Wikipédia (\url{https://fr.wikipedia.org/wiki/H%C3%B4tel_de_Hilbert})
\item Contexte : Ensemble infini au sens de Dedekind --- Wikipédia (en anglais) (\url{https://en.wikipedia.org/wiki/Dedekind-infinite_set})
\item Lecture : Paul Halmos, \textit{Naive Set Theory}
\end{refs}

\bigskip

\begin{problem}[{La récurrence et le principe du plus petit élément --- Lycée, short}]
\textbf{FND-37.} \textbf{Problème.} Le \textit{principe du plus petit élément} dit que toute partie
non vide de l'ensemble des entiers naturels contient un plus petit élément. Démontrez qu'il est
équivalent au principe de récurrence --- déduisez chacun de l'autre.
\end{problem}

\noindent Tout argument du type \og{} prenons un contre-exemple minimal \fg{} et toute récurrence
sont la même démonstration sous des habits différents, et c'est le théorème qui le dit ; la
méthode de descente infinie de Fermat (NT-29, sur la page de théorie des nombres) est le principe
du plus petit élément employé comme une arme. Peano a pris la récurrence pour axiome en 1889, si
bien que chez lui le principe du plus petit élément est un théorème --- mais ce choix est une
convention, non un fait sur les entiers. Un avertissement pour plus tard : restreints à des
formules de complexité logique bornée, les deux schémas cessent d'être interchangeables, et
démêler exactement en quoi ils diffèrent est un vrai sujet, l'étude des fragments de
l'arithmétique.

\begin{refs}
\item Contexte : Principe du bon ordre --- Wikipédia (en anglais) (\url{https://en.wikipedia.org/wiki/Well-ordering_principle})
\item Contexte : Raisonnement par récurrence --- Wikipédia (\url{https://fr.wikipedia.org/wiki/Raisonnement_par_r%C3%A9currence})
\end{refs}

\bigskip


\section*{Licence}

\begin{problem}[{$\mathbb{Q}$ est dénombrable --- Licence}]
\textbf{FND-06.} \textbf{Problème.} On dit qu'un ensemble est \textit{dénombrable} s'il est fini ou en
bijection avec $ \mathbb{N} $.
\begin{enumerate}
\item[(a)] Exhibez des bijections explicites montrant que
$ \mathbb{Z} $ et $ \mathbb{N} \times \mathbb{N} $ sont dénombrables.
\item[(b)] Démontrez que
$ \mathbb{Q} $ est dénombrable.
\item[(c)] Démontrez qu'une réunion dénombrable d'ensembles
dénombrables est dénombrable --- et repérez précisément où votre démonstration utilise l'axiome du
choix.
\end{enumerate}

\end{problem}

\noindent Cantor a établi ces faits dans les années 1870, dans sa correspondance avec
Dedekind, avant que presque personne ne croie que \og{} de même taille \fg{} ait un sens pour des
ensembles infinis. La partie (c) cache une véritable subtilité sur laquelle glissent même des
auteurs soigneux : pour énumérer la réunion, il faut \textit{choisir} simultanément une énumération
de chaque ensemble, et sans une forme de choix cela peut échouer --- dans certains modèles de ZF
sans choix, $\mathbb{R}$ est une réunion dénombrable d'ensembles dénombrables.

\begin{refs}
\item Contexte : Ensemble dénombrable --- Wikipédia (\url{https://fr.wikipedia.org/wiki/Ensemble_d%C3%A9nombrable})
\item Lecture : Herbert Enderton, \textit{Elements of Set Theory}, ch. 6
\item Article : Georg Cantor, "Über eine Eigenschaft des Inbegriffes aller reellen algebraischen Zahlen" (1874), Crelle's Journal
\end{refs}

\bigskip

\begin{problem}[{Diagonalisation : $\mathbb{R}$ n'est pas dénombrable --- Licence}]
\textbf{FND-07.} \textbf{Problème.}
\begin{enumerate}
\item[(a)] Démontrez que $ \mathbb{R} $ n'est pas dénombrable par l'argument
diagonal, en traitant explicitement le fait que certains réels ont deux développements décimaux
($ 0{,}0999\ldots = 0{,}1000\ldots $).
\item[(b)] Démontrez le théorème de Cantor : pour tout ensemble
$ A $, il n'existe pas de surjection de $ A $ sur l'ensemble de ses parties
$ \mathcal{P}(A) $.
\item[(c)] Déduisez-en qu'il existe une hiérarchie infinie strictement croissante de cardinalités
infinies.
\end{enumerate}

\end{problem}

\noindent La première démonstration de non-dénombrabilité de Cantor (1874) utilisait des
intervalles emboîtés ; l'argument diagonal est venu en 1891 et s'est révélé être le passe-partout
du sujet --- le même geste anime le paradoxe de Russell, l'insolubilité du problème de l'arrêt et
le lemme de diagonalisation de Gödel. La diagonalisation comme technique de démonstration est
traitée en détail dans L'art de la démonstration (\url{proofs.html}). L'hostilité de
Kronecker à ces théorèmes, et le serment ultérieur de Hilbert que \og{} nul ne nous chassera du
paradis que Cantor a créé \fg{}, encadrent toute la crise des fondements.

\begin{refs}
\item Contexte : Argument de la diagonale de Cantor --- Wikipédia (\url{https://fr.wikipedia.org/wiki/Argument_de_la_diagonale_de_Cantor})
\item Contexte : Théorème de Cantor --- Wikipédia (\url{https://fr.wikipedia.org/wiki/Th%C3%A9or%C3%A8me_de_Cantor})
\item Lecture : Herbert Enderton, \textit{Elements of Set Theory}, ch. 6
\end{refs}

\bigskip

\begin{problem}[{Le théorème de Cantor--Bernstein --- Licence}]
\textbf{FND-08.} \textbf{Problème.} Supposons qu'il existe des injections $ f\colon A \to B $ et
$ g\colon B \to A $.
\begin{enumerate}
\item[(a)] Démontrez qu'il existe une bijection entre $ A $ et $ B $. Votre démonstration ne doit
pas utiliser l'axiome du choix --- le théorème appartient à ZF.
\item[(b)] Utilisez le théorème, avec des injections explicitement construites dans les deux sens, pour
démontrer que $ \mathbb{R} $ est en bijection avec $ \mathcal{P}(\mathbb{N}) $, l'ensemble
des parties de $ \mathbb{N} $.
\end{enumerate}

\end{problem}

\noindent Cantor a énoncé le théorème mais le déduisait du bon ordre ; des démonstrations
n'utilisant pas le choix ont été trouvées par Dedekind (1887, non publiée de son vivant) et par
Felix Bernstein, élève de Cantor (1897), tandis que la démonstration publiée par Schröder
contenait une erreur. C'est le cheval de trait de l'arithmétique cardinale : pour montrer que deux
ensembles ont la même taille, on n'a jamais besoin d'une bijection astucieuse --- deux injections
faciles suffisent.

\begin{refs}
\item Contexte : Théorème de Cantor--Bernstein --- Wikipédia (\url{https://fr.wikipedia.org/wiki/Th%C3%A9or%C3%A8me_de_Cantor-Bernstein})
\item Lecture : Herbert Enderton, \textit{Elements of Set Theory}, ch. 6
\end{refs}

\bigskip

\begin{problem}[{Les trois visages du choix --- Licence}]
\textbf{FND-09.} \textbf{Problème.} En travaillant dans ZF, démontrez que les énoncés suivants sont équivalents :
\begin{enumerate}
\item[(a)] l'\textit{axiome du choix} --- toute famille d'ensembles non vides admet une fonction de choix ;
\item[(b)] le \textit{lemme de Zorn} --- tout ensemble ordonné non vide dans lequel toute chaîne admet un
majorant contient un élément maximal ;
\item[(c)] le \textit{théorème de Zermelo} --- tout ensemble peut
être bien ordonné. Utilisez ensuite le lemme de Zorn pour démontrer que tout espace vectoriel
possède une base.
\end{enumerate}

\end{problem}

\noindent Zermelo a déduit (iii) de (i) en 1904 et a allumé l'une des controverses les
plus vives de l'histoire des mathématiques ; le lemme qui porte le nom de Zorn (1935) avait été
trouvé plus tôt par Kuratowski (1922). La plaisanterie de Jerry Bona en dit toute la psychologie :
\og{} L'axiome du choix est manifestement vrai, le principe du bon ordre manifestement faux, et qui
peut se prononcer sur le lemme de Zorn ? \fg{} La boucle s'est refermée en 1984, quand Andreas Blass a
démontré que \og{} tout espace vectoriel possède une base \fg{} implique l'axiome du choix sur ZF.

\begin{refs}
\item Contexte : Axiome du choix --- Wikipédia (\url{https://fr.wikipedia.org/wiki/Axiome_du_choix})
\item Contexte : Lemme de Zorn --- Wikipédia (\url{https://fr.wikipedia.org/wiki/Lemme_de_Zorn})
\item Contexte : Théorème de Zermelo (théorème du bon ordre) --- Wikipédia (\url{https://fr.wikipedia.org/wiki/Th%C3%A9or%C3%A8me_de_Zermelo})
\item Lecture : Thomas Jech, \textit{The Axiom of Choice}
\end{refs}

\bigskip

\begin{problem}[{L'arithmétique au-delà du fini --- Licence}]
\textbf{FND-10.} \textbf{Problème.} L'addition, la multiplication et l'exponentiation ordinales se
définissent par récurrence transfinie : $ \alpha + 0 = \alpha $,
$ \alpha + (\beta+1) = (\alpha+\beta) + 1 $, et bornes supérieures aux étapes limites ; de
même pour la multiplication et l'exponentiation.
\begin{enumerate}
\item[(a)] Montrez que $ 1 + \omega = \omega \ne \omega + 1 $ et
que $ 2 \cdot \omega = \omega \ne \omega \cdot 2 $.
\item[(b)] Déterminez, avec démonstration ou
contre-exemple, lesquelles des propriétés suivantes valent pour tous les ordinaux : associativité
de l'addition et de la multiplication ; commutativité de l'addition et de la multiplication ;
distributivité à gauche
$ \alpha \cdot (\beta + \gamma) = \alpha\cdot\beta + \alpha\cdot\gamma $ ; distributivité à
droite $ (\alpha + \beta) \cdot \gamma = \alpha\cdot\gamma + \beta\cdot\gamma $.
\item[(c)] Soit $ \varepsilon_0 = \sup\{\omega,\ \omega^{\omega},\ \omega^{\omega^{\omega}},\
\ldots\} $. Démontrez que $ \omega^{\varepsilon_0} = \varepsilon_0 $ et que
$ \varepsilon_0 $ est le plus petit ordinal ayant cette propriété.
\end{enumerate}

\end{problem}

\noindent Cantor a créé les ordinaux dans ses \textit{Grundlagen} de 1883 pour compter
\og{} au-delà de l'infini \fg{} les étapes d'une construction sur les séries trigonométriques ; von
Neumann en a donné la définition moderne (chaque ordinal est l'ensemble de ses prédécesseurs) en
1923. L'étrange arithmétique asymétrique n'est pas un défaut mais un relevé fidèle de la géométrie
des bons ordres. L'ordinal $ \varepsilon_0 $ de la partie (c) mesure en fait la force exacte de
l'arithmétique de Peano --- c'est la démonstration de cohérence de Gentzen de 1936, et la raison
pour laquelle FND-11 fonctionne.

\begin{refs}
\item Contexte : Arithmétique ordinale --- Wikipédia (\url{https://fr.wikipedia.org/wiki/Op%C3%A9rations_arithm%C3%A9tiques_sur_les_ordinaux})
\item Contexte : Nombre epsilon --- Wikipédia (\url{https://fr.wikipedia.org/wiki/Nombre_epsilon})
\item Contexte : Récurrence transfinie --- Wikipédia (\url{https://fr.wikipedia.org/wiki/R%C3%A9currence_transfinie})
\item Lecture : Herbert Enderton, \textit{Elements of Set Theory}, ch. 7--8
\end{refs}

\bigskip

\begin{problem}[{Les suites de Goodstein --- Licence}]
\textbf{FND-11.} \textbf{Problème.} Écrivez un entier strictement positif en \textit{notation héréditaire en base
n} : exprimez-le en base $ n $, puis exprimez chaque exposant en base $ n $, et ainsi de
suite (par exemple, en base héréditaire 2, $ 26 = 2^{2^2} + 2^{2+1} + 2 $). La suite de
Goodstein de $ m $ commence à $ m $ ; à chaque étape, on réécrit la valeur courante en base
héréditaire $ n $, on remplace chaque occurrence de $ n $ par $ n+1 $, puis on retranche 1
(et on augmente $ n $ de 1 pour l'étape suivante), en partant de $ n = 2 $.
\begin{enumerate}
\item[(a)] Calculez la suite de Goodstein de 3 jusqu'à
son terme, et calculez les quatre premiers termes de la suite de Goodstein de 4.
\item[(b)] Démontrez le
théorème de Goodstein : pour toute valeur initiale $ m $, la suite de Goodstein de $ m $
finit par atteindre 0.
\end{enumerate}

\end{problem}

\noindent Reuben Goodstein a démontré cela en 1944 en utilisant les ordinaux inférieurs à
$ \varepsilon_0 $ --- la croissance de la base est dominée, en un sens précis, par une suite
strictement décroissante d'ordinaux, et FND-10 fournit tout le nécessaire. La suite étonnante est
venue en 1982 : Kirby et Paris ont démontré que le théorème de Goodstein, bien qu'énoncé purement
sur les entiers, est \textit{indémontrable dans l'arithmétique de Peano} --- l'un des premiers exemples
naturels d'incomplétude gödelienne. Leur article démontre la même chose pour la terminaison du jeu
de l'hydre, une bataille d'élagage d'arbres qui vaut la rencontre.

\begin{refs}
\item Contexte : Théorème de Goodstein --- Wikipédia (\url{https://fr.wikipedia.org/wiki/Th%C3%A9or%C3%A8me_de_Goodstein})
\item Contexte : Jeu de l'hydre --- Wikipédia (en anglais) (\url{https://en.wikipedia.org/wiki/Hydra_game})
\item Article : L. Kirby and J. Paris, "Accessible independence results for Peano arithmetic" (1982), Bull. London Math. Soc. 14
\item Article : R. L. Goodstein, "On the restricted ordinal theorem" (1944), J. Symbolic Logic 9
\end{refs}

\bigskip

\begin{problem}[{La compacité en miniature --- Licence}]
\textbf{FND-12.} \textbf{Problème.}
\begin{enumerate}
\item[(a)] Démontrez le théorème de compacité pour la logique
propositionnelle : un ensemble $ \Sigma $ de formules propositionnelles (sur un ensemble
dénombrable de variables) est satisfiable si et seulement si toute partie finie de $ \Sigma $
l'est.
\item[(b)] Déduisez-en le
théorème de De Bruijn--Erdős : pour tout $ k $, un graphe (éventuellement infini) est
$ k $-coloriable si et seulement si chacun de ses sous-graphes finis l'est.
\end{enumerate}

\end{problem}

\noindent La compacité propositionnelle est l'ancêtre combinatoire du théorème de
compacité du premier ordre de FND-13, et l'application au coloriage est due à De Bruijn et Erdős
(1951). Le principe est un authentique fragment du choix : sur ZF, il est équivalent au théorème
de l'idéal premier booléen, strictement plus faible que l'axiome du choix complet. Dans le cas
dénombrable, le lemme de l'infini de Kőnig (1927) --- tout arbre infini à branchement fini possède
une branche infinie --- en est le moteur naturel, et mérite d'être démontré pour lui-même.

\begin{refs}
\item Contexte : Théorème de compacité --- Wikipédia (\url{https://fr.wikipedia.org/wiki/Th%C3%A9or%C3%A8me_de_compacit%C3%A9})
\item Contexte : Lemme de Kőnig --- Wikipédia (\url{https://fr.wikipedia.org/wiki/Lemme_de_K%C5%91nig})
\item Contexte : Théorème de De Bruijn--Erdős (théorie des graphes) --- Wikipédia (\url{https://fr.wikipedia.org/wiki/Th%C3%A9or%C3%A8me_de_De_Bruijn-Erd%C5%91s_(th%C3%A9orie_des_graphes)})
\end{refs}

\bigskip

\begin{problem}[{Pas de descente infinie dans les ordinaux --- Licence, short}]
\textbf{FND-26.} \textbf{Problème.} Démontrez qu'il n'existe pas de suite infinie d'ordinaux
\[ \alpha_0 > \alpha_1 > \alpha_2 > \cdots \]
strictement décroissante à l'infini. Vous pouvez utiliser le fait que tout ensemble non vide
d'ordinaux possède un plus petit élément.
\end{problem}

\noindent Cette seule ligne est tout le moteur de la récurrence transfinie : un argument
qui associe un ordinal décroissant à chaque étape d'un processus doit s'arrêter, si
astronomiquement grands que soient les ordinaux en jeu. C'est pourquoi les suites de Goodstein de
FND-11 redescendent jusqu'à zéro alors même qu'elles grimpent au-delà de toute borne qu'on puisse
énoncer d'avance, et c'est la raison pour laquelle les démonstrations de terminaison en
informatique s'écrivent avec des ordinaux plutôt qu'avec des entiers.

\begin{refs}
\item Contexte : Récurrence transfinie --- Wikipédia (\url{https://fr.wikipedia.org/wiki/R%C3%A9currence_transfinie})
\item Contexte : Nombre ordinal --- Wikipédia (\url{https://fr.wikipedia.org/wiki/Nombre_ordinal})
\end{refs}

\bigskip

\begin{problem}[{Le principe des tiroirs infini, deux fois --- Licence, short}]
\textbf{FND-27.} \textbf{Problème.} Coloriez chaque partie à deux éléments de $\mathbb{N}$ avec l'une d'un
nombre fini de couleurs. Démontrez qu'il existe un ensemble infini $ H \subseteq \mathbb{N} $
dont toutes les parties à deux éléments reçoivent la même couleur.
\end{problem}

\noindent C'est le théorème de Ramsey pour les paires, et la démonstration consiste à
appliquer le principe des tiroirs fini une infinité de fois, puis une fois de plus : on construit
une suite emboîtée d'ensembles infinis, chacun décidant la couleur d'un sommet face à tous les
suivants, puis on applique le principe des tiroirs aux décisions. Frank Ramsey l'a démontré en
1928, à 25 ans et à deux ans de sa mort, comme lemme au sein d'un article sur un problème de
décision pour la logique du premier ordre --- la combinatoire était l'outil, non le but. Remplacez
\og{} paires \fg{} par \og{} sommes finies \fg{} et vous obtenez le théorème de Hindman, dont la force logique
reste inconnue (FND-35).

\begin{refs}
\item Contexte : Théorème de Ramsey --- Wikipédia (\url{https://fr.wikipedia.org/wiki/Th%C3%A9or%C3%A8me_de_Ramsey})
\item Article : F. P. Ramsey, "On a problem of formal logic" (1930), Proc. London Math. Soc. 30
\end{refs}

\bigskip

\begin{problem}[{Ultrafiltres, et des limites qui existent toujours --- Licence}]
\textbf{FND-28.} \textbf{Problème.} Un \textit{filtre} sur un ensemble $ S $ est une famille non vide $ F $ de
parties de $ S $, ne contenant pas $ \varnothing $, stable par intersections finies et par
passage aux sur-ensembles. C'est un \textit{ultrafiltre} si, pour toute partie
$ A \subseteq S $, on a $ A \in F $ ou $ S \setminus A \in F $ ; il est \textit{principal}
s'il est la famille de toutes les parties contenant un point fixé.
\begin{enumerate}
\item[(a)] Démontrez qu'un filtre est un ultrafiltre si et seulement s'il est maximal parmi les filtres
pour l'inclusion.
\item[(b)] À l'aide du lemme de Zorn (FND-09), démontrez que tout filtre s'étend en un ultrafiltre, et
déduisez-en qu'il existe un ultrafiltre non principal sur $\mathbb{N}$ --- c'est-à-dire ne contenant aucune
partie finie.
\item[(c)] Soit $ U $ un ultrafiltre sur $\mathbb{N}$ et $ (x_n) $ une suite réelle bornée. Démontrez qu'il
existe exactement un $ L \in \mathbb{R} $ tel que $ \{ n : |x_n - L| 0 $. Démontrez que cette $ U $-limite est linéaire et monotone en la
suite, et qu'elle coïncide avec $ \lim x_n $ chaque fois que la limite ordinaire existe.
\item[(d)] Démontrez que l'énoncé \og{} tout filtre s'étend en un ultrafiltre \fg{} implique le théorème de
compacité pour la logique propositionnelle (FND-12), et précisez --- références à l'appui --- où se
situe exactement ce principe entre ZF et l'axiome du choix complet.
\end{enumerate}

\end{problem}

\noindent Henri Cartan a introduit les filtres en 1937 pour remplacer les suites en
topologie générale, et Bourbaki les a portés partout ; l'ultrafiltre est l'aboutissement de cette
idée --- un verdict cohérent sur \textit{toute} partie à la fois, une mesure finiment additive à deux
valeurs. La partie (c) est le gain qu'un premier cours montre rarement : un ultrafiltre en main,
toute suite bornée converge, et le prix à payer est qu'on ne peut jamais en écrire un, la
démonstration d'existence étant du pur Zorn. Les ultrafiltres servent ensuite à construire le
compactifié de Stone--Čech, les ultraproduits et le théorème de Łoś, l'analyse non standard de
Robinson, et la démonstration en un paragraphe du théorème de Hindman (FND-35).

\begin{refs}
\item Contexte : Ultrafiltre --- Wikipédia (\url{https://fr.wikipedia.org/wiki/Ultrafiltre}), Filtre (mathématiques) --- Wikipédia (\url{https://fr.wikipedia.org/wiki/Filtre_(math%C3%A9matiques)})
\item Contexte : Théorème de l'idéal premier dans une algèbre de Boole --- Wikipédia (\url{https://fr.wikipedia.org/wiki/Th%C3%A9or%C3%A8me_de_l%27id%C3%A9al_premier_dans_une_alg%C3%A8bre_de_Boole})
\item Lecture : C. C. Chang and H. J. Keisler, \textit{Model Theory}, ch. 4
\item Lecture : Thomas Jech, \textit{The Axiom of Choice}
\end{refs}

\bigskip

\begin{problem}[{Logique du second ordre : la catégoricité payée au prix fort --- Licence}]
\textbf{FND-38.} \textbf{Problème.} En logique du second ordre avec la sémantique \textit{standard}, un
quantificateur $ \forall X $ parcourt \textit{toutes} les parties du domaine --- et non un ensemble
de parties fourni par la structure.
\begin{enumerate}
\item[(a)] Écrivez les axiomes de Peano du second ordre (l'axiome de récurrence est désormais un unique
énoncé
$ \forall X \, [\,(0 \in X \wedge \forall n\,(n \in X \to n+1 \in X)) \to \forall n\, (n \in X)\,] $)
et démontrez le théorème de catégoricité de Dedekind : deux modèles quelconques sont isomorphes.
\item[(b)] Écrivez les axiomes du second ordre d'un corps ordonné complet --- la complétude étant exprimée
par une quantification sur toutes les parties non vides majorées --- et démontrez que deux modèles
quelconques sont isomorphes.
\item[(c)] Déduisez-en que le théorème de compacité (FND-13) et les deux théorèmes de Löwenheim--Skolem
(FND-15) \textit{tombent en défaut} pour la logique du second ordre avec sémantique standard.
Exhibez concrètement l'échec dans chaque cas.
\item[(d)] Démontrez qu'il n'existe aucun système de démonstration correct, complet et effectivement
axiomatisé pour la validité au second ordre. (Combinez (a) avec le premier théorème d'incomplétude
de Gödel, FND-16.) Énoncez ensuite le théorème de Lindström --- la logique du premier ordre est la
plus forte des logiques possédant à la fois la compacité et la propriété de Löwenheim--Skolem
descendante --- et expliquez ce qu'il dit du compromis que vous venez de mesurer.
\end{enumerate}

\end{problem}

\noindent Dedekind a démontré (a) en 1888, avant même que \og{} premier ordre \fg{} et \og{} second
ordre \fg{} ne soient distingués, et pendant une génération la lecture au second ordre a simplement
été ce que \og{} logique \fg{} voulait dire --- les systèmes de Frege et de Russell sont du second ordre. Le
repli sur la logique du premier ordre dans les années 1920 et 1930 fut un échange délibéré, et ce
problème est cet échange mis à plat : la sémantique pleine détermine exactement $ \mathbb{N} $
et $ \mathbb{R} $, ce qui est précisément pourquoi elle ne peut avoir de système de
démonstration utilisable, tandis que la logique du premier ordre a la complétude et la compacité
et ne peut donc rien déterminer d'infini (FND-14). Lindström a démontré en 1969 qu'il n'y a pas de
troisième voie : toute logique possédant les deux propriétés retombe sur le premier ordre. La
sémantique alternative de Henkin (1950) restaure la complétude en restreignant les quantificateurs
du second ordre, et abandonne la catégoricité du même geste. Le verdict de Quine selon lequel la
logique du second ordre est \og{} de la théorie des ensembles déguisée en brebis \fg{} et la réponse de
Shapiro, longue d'un livre, forment la postérité philosophique du théorème.

\begin{refs}
\item Contexte : Logique du second ordre --- Wikipédia (\url{https://fr.wikipedia.org/wiki/Logique_du_second_ordre})
\item Contexte : Théorème de Lindström --- Wikipédia (\url{https://fr.wikipedia.org/wiki/Th%C3%A9or%C3%A8me_de_Lindstr%C3%B6m})
\item Lecture : Herbert Enderton, \textit{A Mathematical Introduction to Logic}, ch. 4
\item Lecture : Stewart Shapiro, \textit{Foundations without Foundationalism: A Case for Second-order Logic} (Oxford, 1991)
\end{refs}

\bigskip

\begin{problem}[{La hiérarchie arithmétique --- Licence}]
\textbf{FND-39.} \textbf{Problème.} Une partie $ A \subseteq \mathbb{N} $ est $ \Sigma^0_n $ si elle est
définissable par une formule à $ n $ blocs de quantificateurs alternés commençant par
$ \exists $, appliqués à une matrice calculable ; $ \Pi^0_n $ est la classe duale, et
$ \Delta^0_n = \Sigma^0_n \cap \Pi^0_n $.
\begin{enumerate}
\item[(a)] Démontrez les faits de clôture élémentaires : chaque $ \Sigma^0_n $ est stable par réunions
finies, intersections finies et quantification existentielle numérique ; $ \Sigma^0_n \cup
\Pi^0_n \subseteq \Delta^0_{n+1} $ ; et $ \Delta^0_1 $ est exactement la classe des ensembles
calculables tandis que $ \Sigma^0_1 $ est exactement celle des ensembles récursivement
énumérables.
\item[(b)] Situez correctement trois ensembles concrets et démontrez vos placements : l'ensemble d'arrêt
$ K $ est $ \Sigma^0_1 $-complet ; $ \mathrm{TOT} = \{ e : \varphi_e \text{ est totale} \} $
est $ \Pi^0_2 $-complet ; et $ \{ e : \varphi_e \text{ a un domaine infini} \} $ est
$ \Pi^0_2 $.
\item[(c)] Démontrez le théorème de hiérarchie : pour tout $ n $, il existe un ensemble
$ \Sigma^0_{n} $ qui n'est pas $ \Pi^0_{n} $, de sorte que la hiérarchie est stricte et ne
s'effondre pas. (Construisez un ensemble $ \Sigma^0_n $ universel et diagonalisez --- le geste de
FND-07, encore une fois.)
\item[(d)] Démontrez le théorème de Post : $ A $ est $ \Sigma^0_{n+1} $ si et seulement si $ A $
est récursivement énumérable relativement au $ n $-ième saut de Turing $ \varnothing^{(n)} $,
et $ \varnothing^{(n)} $ est $ \Sigma^0_n $-complet. Expliquez ce que signifie dire que
complexité de définition et complexité de calcul sont ici une seule et même mesure.
\end{enumerate}

\end{problem}

\noindent Kleene (1943) et Mostowski (1946) sont arrivés à la hiérarchie indépendamment,
et c'est le théorème de Post qui en fait plus qu'un dispositif de comptabilité syntaxique :
compter les alternances de quantificateurs et itérer le problème de l'arrêt se révèlent être deux
noms pour une seule chose. Le gain pratique est que \og{} à quel point est-ce incalculable ? \fg{} devient
une question à réponse numérique --- il suffit de regarder la forme de la définition. La hiérarchie
est aussi la colonne vertébrale de plusieurs entrées plus loin sur cette page : $ \mathrm{ACA}_0 $
en mathématiques à rebours (FND-35, FND-44) est exactement le sous-système clos par saut de
Turing, et l'$ \omega $-cohérence de FND-16 est une hypothèse de saveur $ \Pi^0_2 $ pour la
même raison.

\begin{refs}
\item Contexte : Hiérarchie arithmétique --- Wikipédia (\url{https://fr.wikipedia.org/wiki/Hi%C3%A9rarchie_arithm%C3%A9tique})
\item Contexte : Théorème de Post --- Wikipédia (\url{https://fr.wikipedia.org/wiki/Th%C3%A9or%C3%A8me_de_Post})
\item Lecture : Robert I. Soare, \textit{Recursively Enumerable Sets and Degrees} (Springer, 1987)
\item Lecture : George Boolos, John Burgess, and Richard Jeffrey, \textit{Computability and Logic}
\end{refs}

\bigskip

\begin{problem}[{Réels calculables, et l'égalité que l'on ne peut pas tester --- Licence, short}]
\textbf{FND-40.} \textbf{Problème.} On dit qu'un réel $ \alpha $ est \textit{calculable} si un
algorithme, recevant $ n $, produit un rationnel à $ 2^{-n} $ près de $ \alpha $. Démontrez
que seuls des réels en quantité dénombrable sont calculables --- il existe donc des réels non
calculables --- et démontrez qu'aucun algorithme ne décide, à partir de programmes pour deux réels
calculables, si ces deux réels sont égaux.
\end{problem}

\noindent Turing a introduit les nombres calculables dans l'article de 1936 qui a
introduit les machines de Turing ; les machines étaient le moyen et les nombres étaient le but, ce
qui est l'inverse de ce dont on se souvient d'ordinaire (Borel avait esquissé l'idée en 1912).
L'indécidabilité de l'égalité n'est pas un détail technique : c'est pourquoi toute bibliothèque
d'arithmétique réelle exacte ne compare les nombres qu'à une précision annoncée, et pourquoi les
réels calculables, bien qu'ils forment un corps réel clos, ne se manipulent pas comme les réels
d'un premier cours. Opposez calculabilité et définissabilité --- la probabilité d'arrêt $ \Omega $
de Chaitin est un réel parfaitement spécifié qu'aucun algorithme ne peut approcher --- et vous tenez
le début du sujet que FND-39 mesure.

\begin{refs}
\item Contexte : Nombre réel calculable --- Wikipédia (\url{https://fr.wikipedia.org/wiki/Nombre_r%C3%A9el_calculable})
\item Article : A. M. Turing, "On computable numbers, with an application to the Entscheidungsproblem" (1936), \textit{Proceedings of the London Mathematical Society} 42
\end{refs}

\bigskip


\section*{Master}

\begin{problem}[{Le théorème de compacité --- Master}]
\textbf{FND-13.} \textbf{Problème.}
\begin{enumerate}
\item[(a)] Déduisez le théorème de compacité pour la logique du premier ordre
--- un ensemble d'énoncés admet un modèle si et seulement si toute partie finie en admet un --- du
théorème de complétude de Gödel.
\item[(b)] Utilisez la compacité pour démontrer que la classe des groupes finis n'est pas
axiomatisable au premier ordre : il n'existe aucun ensemble d'énoncés du premier ordre dont les
modèles soient exactement les groupes finis.
\item[(c)] Démontrez le théorème de Löwenheim--Skolem ascendant : une théorie ayant un
modèle infini a des modèles de tout cardinal au moins égal à celui de son langage.
\end{enumerate}

\end{problem}

\noindent Gödel a démontré la complétude et la compacité pour les langages dénombrables
dans sa thèse de 1930 ; Mal'cev a étendu la compacité aux langages quelconques en 1936 et fut le
premier à s'en servir comme d'un outil en algèbre. La compacité est la source vive de la théorie
des modèles : elle fabrique des structures aux propriétés prescrites à partir de pure syntaxe, et
ses conséquences --- la finitude est invisible à la logique du premier ordre, la cardinalité
au-dessus de $ \aleph_0 $ est négociable --- marquent exactement la frontière du pouvoir
expressif du premier ordre.

\begin{refs}
\item Contexte : Théorème de compacité --- Wikipédia (\url{https://fr.wikipedia.org/wiki/Th%C3%A9or%C3%A8me_de_compacit%C3%A9})
\item Contexte : Théorème de complétude de Gödel --- Wikipédia (\url{https://fr.wikipedia.org/wiki/Th%C3%A9or%C3%A8me_de_compl%C3%A9tude_de_G%C3%B6del})
\item Contexte : Théorème de Löwenheim--Skolem --- Wikipédia (\url{https://fr.wikipedia.org/wiki/Th%C3%A9or%C3%A8me_de_L%C3%B6wenheim-Skolem})
\item Lecture : David Marker, \textit{Model Theory: An Introduction}, ch. 2
\end{refs}

\bigskip

\begin{problem}[{Modèles non standard de l'arithmétique --- Master}]
\textbf{FND-14.} \textbf{Problème.} Soit $ \mathrm{Th}(\mathbb{N}) $ l'ensemble de tous les
énoncés du premier ordre vrais dans la structure $ (\mathbb{N};\, 0, 1, +, \cdot, <) $.
\begin{enumerate}
\item[(a)] Utilisez la compacité pour construire un modèle de $ \mathrm{Th}(\mathbb{N}) $ contenant un
élément plus grand que tout entier naturel standard.
\item[(b)] Démontrez que le type d'ordre de tout
modèle non standard dénombrable de $ \mathrm{Th}(\mathbb{N}) $ est $ \mathbb{N} $ suivi d'une
famille dense de copies de $ \mathbb{Z} $ --- autrement dit, $ \mathbb{N} + \mathbb{Z} \cdot
\mathbb{Q} $.
\item[(c)] Dedekind a démontré que l'axiome de récurrence du \textit{second ordre} détermine
$ (\mathbb{N};\, 0, \mathrm{successeur}) $ à isomorphisme près. Expliquez précisément pourquoi
(a) ne le contredit pas : sur quoi le schéma de récurrence du premier ordre quantifie-t-il, et que
manque-t-il ?
\end{enumerate}

\end{problem}

\noindent Skolem a construit des modèles non standard en 1934, montrant qu'aucune liste
de vérités du premier ordre ne détermine les entiers naturels. Le théorème de Tennenbaum (1959)
ajoute une contrainte frappante : dans aucun modèle non standard dénombrable de PA l'addition ou
la multiplication n'est calculable --- ces modèles existent mais ne peuvent jamais être exhibés
chiffre par chiffre. Abraham Robinson a transformé le même phénomène sur $\mathbb{R}$ en atout en 1961 :
l'analyse non standard, où les infiniment petits vivent enfin rigoureusement.

\begin{refs}
\item Contexte : Modèle non standard de l'arithmétique --- Wikipédia (\url{https://fr.wikipedia.org/wiki/Mod%C3%A8le_non_standard_de_l%27arithm%C3%A9tique})
\item Contexte : Théorème de Tennenbaum --- Wikipédia (\url{https://fr.wikipedia.org/wiki/Th%C3%A9or%C3%A8me_de_Tennenbaum})
\item Lecture : Richard Kaye, \textit{Models of Peano Arithmetic}
\end{refs}

\bigskip

\begin{problem}[{Le paradoxe de Skolem --- Master}]
\textbf{FND-15.} \textbf{Problème.}
\begin{enumerate}
\item[(a)] Démontrez le théorème de Löwenheim--Skolem descendant pour un
langage dénombrable : toute structure possède une sous-structure élémentaire dénombrable.
\item[(b)] Supposons ZFC cohérente ; par complétude elle a un modèle, et par (a) elle a un modèle
dénombrable $ M $. À l'intérieur de $ M $, le théorème \og{} $\mathbb{R}$ n'est pas dénombrable \fg{} (FND-07)
est vrai. Résolvez la contradiction apparente en toute rigueur : dites exactement quelle notion
cesse d'être absolue entre $ M $ et l'univers ambiant, exhibez l'énoncé dont le sens se déplace,
et expliquez ce que \og{} dénombrable \fg{} veut dire \textit{dans} $ M $ par opposition à \textit{au sujet
de} $ M $.
\end{enumerate}

\end{problem}

\noindent Thoralf Skolem a présenté cela en 1922 comme un argument montrant que les
fondements ensemblistes sont intrinsèquement relatifs ; la résolution moderne --- la dénombrabilité
d'un ensemble est attestée par une bijection, et cette bijection peut vivre hors du modèle --- est
une leçon définitive de relativisation, et le pain quotidien de quiconque travaille avec des
modèles de la théorie des ensembles. Hilary Putnam a plus tard ranimé le paradoxe en philosophie
sous le nom d'\og{} argument modèle-théorique \fg{}. Le traverser honnêtement est la meilleure préparation
possible au forcing (FND-18).

\begin{refs}
\item Contexte : Paradoxe de Skolem --- Wikipédia (\url{https://fr.wikipedia.org/wiki/Paradoxe_de_Skolem})
\item Contexte : Théorème de Löwenheim--Skolem --- Wikipédia (\url{https://fr.wikipedia.org/wiki/Th%C3%A9or%C3%A8me_de_L%C3%B6wenheim-Skolem})
\item Lecture : Kenneth Kunen, \textit{Set Theory: An Introduction to Independence Proofs}, ch. IV
\end{refs}

\bigskip

\begin{problem}[{Pourquoi \og{} cet énoncé est indémontrable \fg{} exige l'arithmétisation --- Master}]
\textbf{FND-16.} \textbf{Problème.} Soit $ T $ une théorie cohérente, récursivement axiomatisée,
extension de l'arithmétique de Peano, et soit $ \ulcorner \varphi \urcorner $ le
numéral du numéro de Gödel de $ \varphi $.
\begin{enumerate}
\item[(a)] Démontrez le lemme de diagonalisation : pour toute formule
$ \psi(x) $ à une variable libre, il existe un énoncé $ \varphi $ tel que
$ T \vdash \varphi \leftrightarrow \psi(\ulcorner \varphi \urcorner) $.
\item[(b)] Appliquez (a) à
la négation du prédicat de démontrabilité pour obtenir l'énoncé de Gödel $ G $, et démontrez :
$ T $ ne démontre pas $ G $ ; et si $ T $ est $ \omega $-cohérente, $ T $ ne démontre
pas $ \neg G $ non plus.
\item[(c)] Expliquez comment la modification de Rosser (1936) remplace
l'$ \omega $-cohérence par la simple cohérence.
\item[(d)] Démontrez le théorème d'indéfinissabilité de Tarski --- aucune
formule arithmétique $ \mathrm{True}(x) $ ne vérifie $ \mathbb{N} \models \varphi
\leftrightarrow \mathrm{True}(\ulcorner \varphi \urcorner) $ pour tous les énoncés $ \varphi $
--- et servez-vous-en pour expliquer précisément pourquoi l'énoncé de Gödel doit dire \og{} je suis
indémontrable \fg{} plutôt que \og{} je suis faux \fg{}.
\end{enumerate}

\end{problem}

\noindent C'est autant un problème de compréhension qu'un problème de démonstration : la
paraphrase courante \og{} Gödel a formalisé le menteur \fg{} est fautive sur un point essentiel, et les
parties (a) et (d) montrent pourquoi. Un énoncé formel n'a pas de mot \og{} ceci \fg{} ; l'autoréférence
doit être fabriquée par arithmétisation --- en codant la syntaxe par des nombres, de sorte qu'un
énoncé puisse parler de (du code de) lui-même via la fonction de substitution. La démontrabilité
est arithmétiquement définissable ; la vérité, d'après (d), ne l'est pas. Cette asymétrie est tout
le moteur de l'article de Gödel de 1931, et la raison pour laquelle l'incomplétude est un théorème
et non un paradoxe.

\begin{refs}
\item Contexte : Théorèmes d'incomplétude de Gödel --- Wikipédia (\url{https://fr.wikipedia.org/wiki/Th%C3%A9or%C3%A8mes_d%27incompl%C3%A9tude_de_G%C3%B6del})
\item Contexte : Lemme de diagonalisation --- Wikipédia (en anglais) (\url{https://en.wikipedia.org/wiki/Diagonal_lemma})
\item Contexte : Théorème de Tarski (indéfinissabilité de la vérité) --- Wikipédia (\url{https://fr.wikipedia.org/wiki/Th%C3%A9or%C3%A8me_de_Tarski})
\item Lecture : Peter Smith, \textit{An Introduction to Gödel's Theorems}
\end{refs}

\bigskip

\begin{problem}[{Le théorème de Löb --- Master}]
\textbf{FND-17.} \textbf{Problème.} Soit $ T $ comme en FND-16, et soit
$ \mathrm{Prov}(x) $ un prédicat de démontrabilité standard vérifiant les
conditions de dérivabilité de Hilbert--Bernays--Löb : (D1) si $ T \vdash \varphi $ alors
$ T \vdash \mathrm{Prov}(\ulcorner \varphi \urcorner) $ ;
(D2) $ T \vdash \mathrm{Prov}(\ulcorner \varphi \to \psi \urcorner) \to
(\mathrm{Prov}(\ulcorner \varphi \urcorner) \to \mathrm{Prov}(\ulcorner \psi \urcorner)) $;
(D3) $ T \vdash \mathrm{Prov}(\ulcorner \varphi \urcorner) \to
\mathrm{Prov}(\ulcorner \mathrm{Prov}(\ulcorner \varphi \urcorner) \urcorner) $.
\begin{enumerate}
\item[(a)] Démontrez le théorème de Löb : si $ T \vdash \mathrm{Prov}(\ulcorner \varphi \urcorner) \to
\varphi $, alors $ T \vdash \varphi $.
\item[(b)] Déduisez-en le second théorème d'incomplétude de Gödel : si
$ T $ est cohérente, $ T $ ne démontre pas $ \mathrm{Con}(T) $.
\item[(c)] Déduisez-en que l'énoncé
de Henkin --- le point fixe affirmant \og{} je suis démontrable \fg{} --- est démontrable dans $ T $.
\end{enumerate}

\end{problem}

\noindent En 1952, Leon Henkin a posé une question d'apparence innocente : l'énoncé de
Gödel affirme sa propre indémontrabilité et est indémontrable --- qu'en est-il de l'énoncé affirmant
sa propre démontrabilité ? La réponse de Martin Löb en 1955 est l'un des vrais chocs de la
logique : une théorie démontre \og{} la démontrabilité de $\varphi$ implique $\varphi$ \fg{} seulement quand elle démontre
déjà $\varphi$ purement et simplement. La logique modale GL, la boîte se lisant Prov, capture exactement
les principes modaux démontrables de l'arithmétique --- c'est le théorème de complétude arithmétique
de Solovay (1976), le couronnement de la logique de la démontrabilité.

\begin{refs}
\item Contexte : Théorème de Löb --- Wikipédia (\url{https://fr.wikipedia.org/wiki/Th%C3%A9or%C3%A8me_de_L%C3%B6b})
\item Article : M. H. Löb, "Solution of a problem of Leon Henkin" (1955), J. Symbolic Logic 20
\item Lecture : George Boolos, \textit{The Logic of Provability}
\end{refs}

\bigskip

\begin{problem}[{Comment l'hypothèse du continu s'est révélée indépendante --- Master}]
\textbf{FND-18.} \textbf{Problème.} L'hypothèse du continu (HC) affirme que toute partie infinie
de $\mathbb{R}$ est en bijection avec $\mathbb{N}$ ou avec $\mathbb{R}$, c'est-à-dire $ 2^{\aleph_0} = \aleph_1 $.
\begin{enumerate}
\item[(a)] Définissez la hiérarchie constructible $ L $ de Gödel (à chaque étape, on ne prend que les
parties définissables avec paramètres sur l'étape précédente) et, en admettant le lemme de
condensation, démontrez que l'hypothèse du continu généralisée est vraie dans $ L $ ; concluez
$ \mathrm{Con}(\mathrm{ZF}) \Rightarrow \mathrm{Con}(\mathrm{ZFC} + \mathrm{HCG}) $.
\item[(b)] Énoncez précisément ce que réalise la construction de forcing de Cohen ---
$ \mathrm{Con}(\mathrm{ZF}) \Rightarrow \mathrm{Con}(\mathrm{ZFC} + \neg\mathrm{HC}) $ ---
et expliquez le rôle du modèle de base transitif dénombrable (FND-15) et pourquoi l'ajout de
$ \aleph_2 $ réels de Cohen exige de démontrer que les cardinaux sont préservés (la condition
de chaîne dénombrable).
\item[(c)] Dites exactement ce que les deux moitiés ensemble établissent, et n'établissent pas, quant à
la vérité de HC.
\end{enumerate}

\end{problem}

\noindent HC était le premier des problèmes de Hilbert de 1900, et Cantor l'a poursuivie
jusqu'aux limites de sa santé. Gödel a construit $ L $ et démontré la moitié de l'indépendance
en 1938 ; Paul Cohen, analyste venu de l'extérieur de la logique, a inventé le forcing et l'a
achevée en 1963, ce qui lui a valu la seule médaille Fields jamais décernée pour la logique
(1966). Savoir si l'indépendance clôt la question ou ne fait que la déplacer reste un débat vivant
--- voyez FND-21. Les traitements approfondis de problèmes de ce rang se trouvent sur la page
Problèmes ouverts (\url{open-problems.html}).

\begin{refs}
\item Contexte : Hypothèse du continu --- Wikipédia (\url{https://fr.wikipedia.org/wiki/Hypoth%C3%A8se_du_continu})
\item Contexte : Univers constructible --- Wikipédia (\url{https://fr.wikipedia.org/wiki/Univers_constructible})
\item Contexte : Forcing --- Wikipédia (\url{https://fr.wikipedia.org/wiki/Forcing})
\item Article : Paul Cohen, "The independence of the continuum hypothesis" (1963), Proc. Nat. Acad. Sci. USA 50
\end{refs}

\bigskip

\begin{problem}[{Un îlot de décidabilité : l'arithmétique de Presburger --- Master}]
\textbf{FND-19.} \textbf{Problème.} L'arithmétique de Presburger est la théorie du premier ordre de
$ (\mathbb{N};\, 0, 1, +, <) $ --- l'addition, mais pas la multiplication.
\begin{enumerate}
\item[(a)] Démontrez que cette
théorie admet l'élimination des quantificateurs dès que l'on enrichit le langage des prédicats
$ x \equiv r \pmod{n} $ pour chaque $ n \ge 2 $ et $ 0 \le r < n $.
\item[(b)] Concluez que
l'arithmétique de Presburger est complète et décidable.
\item[(c)] Combinez (b) avec le premier
théorème d'incomplétude de Gödel (FND-16) pour démontrer que la multiplication n'est pas
définissable par une formule du premier ordre dans $ (\mathbb{N};\, 0, 1, +, <) $.
\end{enumerate}

\end{problem}

\noindent Mojżesz Presburger a démontré la décidabilité en 1929, alors qu'il était
étudiant de maîtrise au séminaire de Tarski à Varsovie --- à la légère déception de Tarski,
dit-on, puisqu'il ne s'agissait \og{} que \fg{} d'un problème de maîtrise ; Presburger fut plus tard
assassiné pendant la Shoah. La théorie se tient exactement à la frontière : une multiplication
la sépare de l'indécidabilité. Fischer et Rabin ont démontré en 1974 que toute procédure de
décision exige un temps doublement exponentiel dans le pire des cas, et pourtant des solveurs
Presburger tournent chaque jour au sein des compilateurs et des vérificateurs de programmes
fondés sur SMT.

\begin{refs}
\item Contexte : Arithmétique de Presburger --- Wikipédia (\url{https://fr.wikipedia.org/wiki/Arithm%C3%A9tique_de_Presburger})
\item Lecture : George Boolos, John Burgess, and Richard Jeffrey, \textit{Computability and Logic}
\item Article : M. J. Fischer and M. O. Rabin, "Super-exponential complexity of Presburger arithmetic" (1974), SIAM--AMS Proceedings 7
\end{refs}

\bigskip

\begin{problem}[{Le principe de réflexion --- Master, short}]
\textbf{FND-29.} \textbf{Problème.} Soit $ V_\alpha $ la hiérarchie cumulative
($ V_0 = \varnothing $, $ V_{\alpha+1} = \mathcal{P}(V_\alpha) $, réunions aux limites) et soit
$ \varphi(x_1, \ldots, x_n) $ une formule quelconque du langage de la théorie des ensembles.
Démontrez dans ZF que, pour tout ordinal $ \beta $, il existe un ordinal $ \alpha > \beta $ tel
que, pour tous $ a_1, \ldots, a_n \in V_\alpha $, $ \varphi(a_1, \ldots, a_n) $ est vraie si et
seulement si sa relativisée $ \varphi^{V_\alpha}(a_1, \ldots, a_n) $ l'est.
\end{problem}

\noindent Montague et Lévy ont isolé ce schéma vers 1960, et il dit quelque chose de
troublant : aucun énoncé isolé ne peut distinguer l'univers entier des ensembles de l'un de ses
propres segments initiaux. Combiné au second théorème d'incomplétude de Gödel, il démontre que ZF
n'est pas finiment axiomatisable, et combiné à Löwenheim--Skolem (FND-15), il produit les modèles
transitifs dénombrables au-dessus desquels on effectue le forcing (FND-32).

\begin{refs}
\item Contexte : Principe de réflexion --- Wikipédia (en anglais) (\url{https://en.wikipedia.org/wiki/Reflection_principle})
\item Contexte : Univers de von Neumann --- Wikipédia (\url{https://fr.wikipedia.org/wiki/Univers_de_von_Neumann})
\item Lecture : Kenneth Kunen, \textit{Set Theory: An Introduction to Independence Proofs}, ch. IV
\end{refs}

\bigskip

\begin{problem}[{Le deuxième niveau de la hiérarchie de Borel n'est pas le premier --- Master, short}]
\textbf{FND-30.} \textbf{Problème.} Démontrez que $\mathbb{Q}$ est une partie $ F_\sigma $ de $\mathbb{R}$ --- une
réunion dénombrable de fermés --- mais n'est pas un $ G_\delta $, c'est-à-dire une intersection
dénombrable d'ouverts. Concluez que la hiérarchie de Borel sur $\mathbb{R}$ ne s'effondre pas à son deuxième
niveau.
\end{problem}

\noindent C'est le théorème de Baire qui fait le travail, et le corollaire mérite d'être
énoncé pour lui-même : il n'existe aucune fonction $ f : \mathbb{R} \to \mathbb{R} $ continue
exactement aux points rationnels, car l'ensemble de continuité d'une fonction quelconque est un
$ G_\delta $. Baire a démontré son théorème dans sa thèse de 1899 ; Lebesgue a montré en 1905
que la hiérarchie reste stricte à tout niveau dénombrable, ce qui fait de la théorie descriptive
des ensembles une discipline plutôt qu'un exercice de classification.

\begin{refs}
\item Contexte : Hiérarchie de Borel --- Wikipédia (\url{https://fr.wikipedia.org/wiki/Hi%C3%A9rarchie_de_Borel})
\item Contexte : Théorème de Baire --- Wikipédia (\url{https://fr.wikipedia.org/wiki/Th%C3%A9or%C3%A8me_de_Baire})
\item Lecture : A. S. Kechris, \textit{Classical Descriptive Set Theory}, ch. II
\end{refs}

\bigskip

\begin{problem}[{Le dixième problème de Hilbert et le théorème MRDP --- Master}]
\textbf{FND-31.} \textbf{Problème.} On dit que $ S \subseteq \mathbb{N}^k $ est \textit{diophantien} s'il existe un
polynôme $ P $ à coefficients entiers tel que
\[ \bar{a} \in S \iff \exists\, y_1, \ldots, y_m \in \mathbb{N} : \;
P(\bar{a}, y_1, \ldots, y_m) = 0 . \]
\begin{enumerate}
\item[(a)] Démontrez que les ensembles diophantiens sont stables par réunion, intersection et
quantification existentielle sur $\mathbb{N}$, et qu'un système fini d'équations diophantiennes peut toujours
être remplacé par une seule équation.
\item[(b)] Démontrez que tout ensemble diophantien est récursivement énumérable, en décrivant la
recherche qui l'énumère.
\item[(c)] Le théorème MRDP énonce la réciproque : tout ensemble récursivement énumérable est
diophantien. En l'admettant, démontrez que le dixième problème de Hilbert est indécidable --- il
n'existe aucun algorithme qui décide si une équation polynomiale donnée à coefficients entiers a
une solution en nombres entiers --- par réduction depuis le problème de l'arrêt.
\item[(d)] Déduisez-en qu'il existe un polynôme à coefficients entiers dont l'ensemble des valeurs
positives, quand ses variables parcourent $\mathbb{N}$, est exactement l'ensemble des nombres premiers.
Expliquez soigneusement pourquoi ce n'est pas une formule donnant les nombres premiers.
\end{enumerate}

\end{problem}

\noindent Hilbert a demandé en 1900 une procédure pour décider n'importe quelle équation
diophantienne ; il supposait qu'il en existait une, la notion d'\og{} aucun algorithme \fg{} ne devant
apparaître que trente ans plus tard. Martin Davis, Hilary Putnam et Julia Robinson ont ramené le
problème en 1961 à la recherche d'une seule relation diophantienne à croissance exponentielle ---
l'hypothèse de Robinson, qu'elle a poursuivie pendant deux décennies --- et en janvier 1970, Iouri
Matiiassevitch, âgé de 22 ans, l'a fournie à l'aide des nombres de Fibonacci, dont la croissance
est exactement celle qu'il fallait. Le polynôme générateur de nombres premiers de la partie (d) a
été explicité par Jones, Sato, Wada et Wiens en 1976 : 26 variables, une par lettre de l'alphabet,
de degré 25.

\begin{refs}
\item Contexte : Dixième problème de Hilbert --- Wikipédia (\url{https://fr.wikipedia.org/wiki/Dixi%C3%A8me_probl%C3%A8me_de_Hilbert})
\item Contexte : Ensemble diophantien --- Wikipédia (\url{https://fr.wikipedia.org/wiki/Diophantien})
\item Lecture : Yuri Matiyasevich, \textit{Hilbert's Tenth Problem} (MIT Press, 1993)
\item Article : Martin Davis, "Hilbert's tenth problem is unsolvable" (1973), \textit{American Mathematical Monthly} 80
\end{refs}

\bigskip

\begin{problem}[{Construire une extension générique à la main --- Master}]
\textbf{FND-32.} \textbf{Problème.} FND-18 vous demandait de dire ce que réalise le forcing. Celui-ci vous demande
de le construire. Soit $ M $ un modèle transitif dénombrable de ZFC, et, à l'intérieur de
$ M $, soit $ \mathbb{P} $ l'ensemble des fonctions partielles finies de
$ \omega_2 \times \omega $ dans $ \{0,1\} $, ordonné par inclusion inverse (plus fort veut dire
plus d'information).
\begin{enumerate}
\item[(a)] Démontrez que, $ M $ étant dénombrable, il existe un filtre $ G \subseteq \mathbb{P} $
rencontrant toute partie dense de $ \mathbb{P} $ appartenant à $ M $ --- et démontrez qu'aucun
tel $ G $ ne peut lui-même appartenir à $ M $.
\item[(b)] Définissez les $ \mathbb{P} $-noms et l'évaluation $ \tau \mapsto \tau_G $, et démontrez
que $ M[G] = \{ \tau_G : \tau \in M^{\mathbb{P}} \} $ est transitif, contient $ M $ et
$ G $, et a les mêmes ordinaux que $ M $. Vérifiez l'extensionnalité, la paire et la réunion
dans $ M[G] $, et décrivez comment se traitent les axiomes restants.
\item[(c)] Énoncez la relation de forcing $ p \Vdash \varphi $ et démontrez le lemme de vérité :
$ M[G] \models \varphi $ exactement lorsqu'un $ p \in G $ force $ \varphi $. Expliquez
pourquoi la relation doit être définissable à l'intérieur de $ M $ pour que toute la méthode
fonctionne.
\item[(d)] Démontrez que $ \mathbb{P} $ vérifie la condition de chaîne dénombrable (utilisez le lemme
du $ \Delta $-système), déduisez-en que $ M $ et $ M[G] $ ont les mêmes cardinaux, et
concluez que $ M[G] \models 2^{\aleph_0} \ge \aleph_2 $, de sorte que HC y est fausse.
\end{enumerate}

\end{problem}

\noindent Paul Cohen a inventé cette machinerie en 1963 et elle reste la technique la plus
puissante de la théorie des ensembles ; presque tous les résultats d'indépendance de cette page en
sont une application. Si elle mérite un problème à elle seule, c'est que le forcing est
notoirement difficile à \textit{comprendre} plutôt qu'à employer : le $ G $ générique n'existe pas
dans $ M $, et pourtant $ M $ peut raisonner sur ce qui serait vrai s'il existait, et c'est le
lemme de vérité qui liquide cette circularité. La reformulation à valeurs booléennes de Scott,
Solovay et Vopěnka a été inventée précisément pour rendre le même contenu inévitable ; lire les
deux présentations est le chemin le plus rapide vers l'idée.

\begin{refs}
\item Contexte : Forcing --- Wikipédia (\url{https://fr.wikipedia.org/wiki/Forcing})
\item Contexte : Condition de chaîne dénombrable --- Wikipédia (\url{https://fr.wikipedia.org/wiki/Condition_de_cha%C3%AEne_d%C3%A9nombrable})
\item Lecture : Kenneth Kunen, \textit{Set Theory: An Introduction to Independence Proofs}, ch. VII
\item Article : Timothy Y. Chow, "A beginner's guide to forcing" (2008), arXiv:0712.1320 (\url{https://arxiv.org/abs/0712.1320})
\end{refs}

\bigskip

\begin{problem}[{Cofinalité et théorème de König --- Master}]
\textbf{FND-41.} \textbf{Problème.} La \textit{cofinalité} $ \operatorname{cf}(\alpha) $ d'un ordinal limite
$ \alpha $ est le plus petit type d'ordre d'une partie de $ \alpha $ non bornée dans
$ \alpha $. Un cardinal $ \kappa $ est \textit{régulier} si
$ \operatorname{cf}(\kappa) = \kappa $, et \textit{singulier} sinon.
\begin{enumerate}
\item[(a)] Démontrez que $ \operatorname{cf}(\alpha) $ est toujours un cardinal régulier. Calculez
$ \operatorname{cf}(\aleph_\omega) $ et $ \operatorname{cf}(\aleph_{\omega_1}) $, et
démontrez (en utilisant le choix) que tout cardinal successeur $ \aleph_{\alpha+1} $ est
régulier.
\item[(b)] Démontrez le théorème de König : si $ \kappa_i < \lambda_i $ pour tout $ i $ d'un
ensemble d'indices $ I $, alors
\[ \sum_{i \in I} \kappa_i \;<\; \prod_{i \in I} \lambda_i . \]
Notez où l'axiome du choix est utilisé, et notez que c'est l'une des très rares inégalités
strictes de l'arithmétique cardinale.
\item[(c)] Déduisez-en que $ \kappa < \kappa^{\operatorname{cf}(\kappa)} $ pour tout $ \kappa $
infini, d'où $ \operatorname{cf}(2^{\aleph_0}) > \aleph_0 $, d'où
$ 2^{\aleph_0} \ne \aleph_\omega $. Démontrez plus généralement que
$ \operatorname{cf}(2^{\kappa}) > \kappa $.
\item[(d)] Énoncez le théorème d'Easton (1970) : hormis la monotonie et la contrainte de (c), ZFC ne
démontre absolument rien sur la fonction continu aux cardinaux \textit{réguliers}. Expliquez
précisément pourquoi la méthode d'Easton ne dit rien des cardinaux singuliers, et indiquez où
vivent les vrais théorèmes de ce côté-là --- le théorème de Silver et la théorie pcf de Shelah
(FND-34).
\end{enumerate}

\end{problem}

\noindent Julius König a annoncé au congrès de Heidelberg de 1904 que le continu n'est pas
un aleph --- c'est-à-dire que $\mathbb{R}$ ne peut être bien ordonné --- et Cantor était dans la salle. En un
jour, Zermelo avait trouvé l'erreur. Ce qui a survécu au naufrage est le théorème de la partie
(b), et il vaut la peine d'apprécier à quel point c'est peu et beaucoup : pour l'essentiel, la
\textit{seule} restriction que ZFC impose à la taille du continu est que sa cofinalité soit non
dénombrable. Tout le reste dans cette direction appartient à Easton, et c'est une licence plutôt
qu'une contrainte : $ 2^{\aleph_0} $ peut valoir, de façon cohérente, $ \aleph_1 $,
$ \aleph_2 $, $ \aleph_{17} $ ou $ \aleph_{\omega+1} $. La découverte que les cardinaux
singuliers ne jouissent \textit{pas} de la même liberté --- que ZFC y démontre de vrais théorèmes ---
est la surprise de Silver en 1974 et de la théorie pcf de Shelah, et c'est pourquoi FND-34 est un
problème vivant et non une affaire réglée.

\begin{refs}
\item Contexte : Cofinalité --- Wikipédia (\url{https://fr.wikipedia.org/wiki/Cofinalit%C3%A9})
\item Contexte : Théorème de König (théorie des ensembles) --- Wikipédia (\url{https://fr.wikipedia.org/wiki/Th%C3%A9or%C3%A8me_de_K%C3%B6nig_(th%C3%A9orie_des_ensembles)})
\item Lecture : Thomas Jech, \textit{Set Theory} (Third Millennium Edition), ch. 3 and 5
\item Lecture : Kenneth Kunen, \textit{Set Theory: An Introduction to Independence Proofs}
\end{refs}

\bigskip

\begin{problem}[{Le problème de Souslin et l'axiome de Martin --- Master}]
\textbf{FND-42.} \textbf{Problème.} Cantor a caractérisé $\mathbb{R}$ à isomorphisme d'ordre près comme l'unique ordre total
dense et complet sans extrémités possédant une partie dense dénombrable. Souslin a demandé en 1920
si \og{} posséder une partie dense dénombrable \fg{} peut être affaibli en la \textit{condition de chaîne
dénombrable} : toute famille d'intervalles ouverts non vides deux à deux disjoints est
dénombrable. Une \textit{droite de Souslin} est un contre-exemple ; l'\textit{hypothèse de Souslin}
affirme qu'il n'en existe pas.
\begin{enumerate}
\item[(a)] Démontrez la caractérisation de Cantor, afin que le contenu exact de la question de Souslin
soit clair.
\item[(b)] Un \textit{arbre de Souslin} est un arbre de hauteur $ \omega_1 $ sans chaîne non dénombrable
ni antichaîne non dénombrable. Démontrez qu'il existe une droite de Souslin si et seulement s'il
existe un arbre de Souslin.
\item[(c)] Énoncez l'axiome de Martin $ \mathrm{MA}(\kappa) $ : pour tout ordre partiel vérifiant la
condition de chaîne dénombrable et toute famille d'au plus $ \kappa $ parties denses, il existe
un filtre les rencontrant toutes. Démontrez $ \mathrm{MA}(\aleph_0) $ directement dans ZFC
(c'est le lemme de Rasiowa--Sikorski), démontrez que $ \mathrm{MA}(2^{\aleph_0}) $ est faux, et
démontrez que $ \mathrm{MA}(\aleph_1) $ entraîne qu'il n'existe pas d'arbre de Souslin.
\item[(d)] Énoncez le théorème de Jensen selon lequel le principe diamant $ \diamondsuit $ entraîne
l'existence d'un arbre de Souslin, et que $ \diamondsuit $ est vrai dans le $ L $ de Gödel
(FND-18). Assemblez les pièces en une démonstration complète de l'indépendance de l'hypothèse de
Souslin vis-à-vis de ZFC, en nommant quelle moitié provient de quelle construction.
\end{enumerate}

\end{problem}

\noindent Mikhaïl Souslin l'a posé comme problème 3 dans le premier volume de
\textit{Fundamenta Mathematicae} en 1920 ; il a paru après sa mort à vingt-cinq ans, et est resté
sans réponse près de cinquante ans. Les deux moitiés de la réponse sont arrivées en quatre ans :
Jech (1967) et Tennenbaum (1968) ont forcé l'existence de droites de Souslin, Jensen a extrait
$ \diamondsuit $ de la structure fine de $ L $ et les a obtenues gratuitement là, et Solovay
et Tennenbaum (1971) ont inventé le forcing itéré pour tuer tous les arbres de Souslin d'un coup
--- la construction dont Martin et Solovay (1970) ont ensuite distillé l'axiome de Martin comme
principe autonome. C'est la vraie raison pour laquelle ce problème vaut la peine d'être fait : MA
est le substitut de HC pour le théoricien des ensembles au travail, un axiome compatible avec un
$ 2^{\aleph_0} $ grand qui tranche pourtant quantité de questions combinatoires, et le problème
de Souslin est son lieu de naissance.

\begin{refs}
\item Contexte : Problème de Souslin --- Wikipédia (\url{https://fr.wikipedia.org/wiki/Probl%C3%A8me_de_Souslin})
\item Contexte : Axiome de Martin --- Wikipédia (\url{https://fr.wikipedia.org/wiki/Axiome_de_Martin}), Principe diamant --- Wikipédia (en anglais) (\url{https://en.wikipedia.org/wiki/Diamond_principle})
\item Lecture : Kenneth Kunen, \textit{Set Theory: An Introduction to Independence Proofs}, ch. II
\item Article : R. M. Solovay and S. Tennenbaum, "Iterated Cohen extensions and Souslin's problem", \textit{Annals of Mathematics} 94 (1971), 201--245
\end{refs}

\bigskip

\begin{problem}[{Interpolation et définissabilité --- Master, short}]
\textbf{FND-43.} \textbf{Problème.} Énoncez précisément le théorème d'interpolation de Craig pour
la logique du premier ordre, et servez-vous-en pour démontrer le théorème de définissabilité de
Beth : un symbole de relation dont l'interprétation est uniquement déterminée, dans tout modèle
d'une théorie $ T $, par le reste de la structure, est déjà explicitement définissable, dans
$ T $, par une formule du langage plus petit.
\end{problem}

\noindent Beth a répondu en 1953 à une question soulevée par Padoa au tournant du siècle
--- la méthode de Padoa montre qu'une notion primitive n'est \textit{pas} définissable en exhibant
deux modèles s'accordant par ailleurs, et le théorème de Beth dit que cette méthode est complète,
que définissabilité implicite et explicite coïncident. Craig a démontré son théorème
d'interpolation en 1957 à partir du théorème de Herbrand--Gentzen, et la déduction de Beth à partir
de Craig tient en trois lignes une fois la bonne théorie écrite, ce qui vaut la peine de la faire
à la main. Les deux résultats ont connu une seconde carrière en informatique : les interpolants de
Craig extraits de démonstrations par réfutation sont un moteur standard de la vérification de
modèles fondée sur SAT.

\begin{refs}
\item Contexte : Théorème d'interpolation de Craig --- Wikipédia (\url{https://fr.wikipedia.org/wiki/Th%C3%A9or%C3%A8me_d%27interpolation_de_Craig})
\item Contexte : Théorème de définissabilité de Beth --- Wikipédia (\url{https://fr.wikipedia.org/wiki/Th%C3%A9or%C3%A8me_de_d%C3%A9finissabilit%C3%A9_de_Beth})
\item Lecture : C. C. Chang and H. J. Keisler, \textit{Model Theory}
\item Lecture : Wilfrid Hodges, \textit{A Shorter Model Theory} (Cambridge, 1997)
\end{refs}

\bigskip


\section*{Recherche}

\begin{problem}[{La hiérarchie des grands cardinaux : une ligne ou plusieurs ? --- Recherche}]
\textbf{FND-20.} \textbf{Problème.} Un cardinal $ \kappa $ est \textit{inaccessible} s'il est
non dénombrable, régulier (non borne supérieure de moins de $ \kappa $ cardinaux plus petits),
et limite fort ($ \lambda < \kappa $ entraîne $ 2^{\lambda} < \kappa $).
\begin{enumerate}
\item[(a)] Démontrez
que si $ \kappa $ est inaccessible, alors $ V_{\kappa} $, l'étage de la hiérarchie cumulative
en $ \kappa $, est un modèle de ZFC ; concluez, via le second théorème d'incomplétude de Gödel,
que ZFC, si elle est cohérente, ne peut démontrer l'existence de cardinaux inaccessibles.
\item[(b)] La question
de recherche : les axiomes de grands cardinaux connus --- inaccessible, mesurable, de Woodin,
supercompact, et des dizaines d'autres --- semblent, sans exception, linéairement ordonnés par force
de cohérence. Existe-t-il un théorème précis expliquant cette linéarité observée, ou un couple
naturel d'hypothèses de grands cardinaux de forces de cohérence incomparables ?
\textit{Statut : la linéarité est un phénomène empirique, non un théorème démontré ; formuler même la
question avec précision fait partie du problème, et celui-ci est ouvert.}
\end{enumerate}

\end{problem}

\noindent Les grands cardinaux descendent de Hausdorff et d'Ulam au début du vingtième
siècle et sont devenus, par les travaux de Scott, Solovay, Martin, Woodin et bien d'autres,
l'étalon standard de la force des énoncés mathématiques : on calibre un énoncé indémontrable en
demandant de quel grand cardinal il a besoin. Que tout énoncé naturel se soit jusqu'ici rangé
quelque part sur une seule ligne est l'une des régularités inexpliquées les plus profondes des
mathématiques --- le programme de Gödel pour trancher des questions comme HC par de nouveaux axiomes
repose dessus.

\begin{refs}
\item Contexte : Grand cardinal --- Wikipédia (\url{https://fr.wikipedia.org/wiki/Grand_cardinal})
\item Contexte : Liste d'énoncés indécidables dans ZFC --- Wikipédia (\url{https://fr.wikipedia.org/wiki/Liste_d%27%C3%A9nonc%C3%A9s_ind%C3%A9cidables_dans_ZFC})
\item Lecture : Akihiro Kanamori, \textit{The Higher Infinite}
\end{refs}

\bigskip

\begin{problem}[{Ultimate L --- Recherche}]
\textbf{FND-21.} \textbf{Problème.} Le $ L $ de Gödel (FND-18) est un univers intérieur canonique et ordonné,
mais il est incompatible avec les cardinaux mesurables (Scott, 1961). Le \textit{programme Ultimate-L}
de Woodin demande s'il existe une version \og{} ultime \fg{} de $ L $ : un modèle intérieur canonique
compatible avec tous les grands cardinaux, tel que l'axiome \og{} V = Ultimate L \fg{} tranche
l'hypothèse du continu (par l'affirmative), soit insensible au forcing, et fournisse une théorie
complète des ensembles héréditairement non dénombrables modulo les grands cardinaux.
\begin{enumerate}
\item[(a)] À partir de
l'article de synthèse cité ci-dessous, énoncez précisément le théorème de dichotomie HOD de
Woodin : s'il existe un cardinal extensible, alors le modèle intérieur HOD des ensembles
héréditairement définissables par des ordinaux est soit \og{} proche \fg{} de $ V $, soit \og{} loin \fg{} de
$ V $ en un sens nettement défini.
\item[(b)] Énoncez la conjecture
Ultimate-L et expliquez ce qu'une réponse positive signifierait pour HC. \textit{Statut : la conjecture
Ultimate-L et la conjecture HOD qui l'accompagne sont ouvertes en 2026 ; le programme est actif,
avec des progrès récents sur la dichotomie HOD (par exemple l'affaiblissement de l'hypothèse vers
la compacité forte), et il est contesté --- d'autres théoriciens des ensembles s'attendent à ce que
l'emportent les axiomes de forcing, qui réfutent HC.}
\end{enumerate}

\end{problem}

\noindent C'est la réponse moderne la plus tranchante à \og{} HC a-t-elle jamais été vraiment
tranchée ? \fg{} Woodin lui-même a autrefois argumenté, à partir de sa $\Omega$-logique, que HC est fausse,
puis a changé de cap à mesure qu'émergeait la structure fine d'Ultimate L --- un rare revirement
public au sommet d'un domaine. Le programme serait, s'il aboutissait, ce qui se rapproche le plus
d'un achèvement du programme de Gödel de 1947 pour de nouveaux axiomes.

\begin{refs}
\item Contexte : W. Hugh Woodin --- Wikipédia (\url{https://fr.wikipedia.org/wiki/W._Hugh_Woodin})
\item Contexte : Ensemble définissable par des ordinaux --- Wikipédia (en anglais) (\url{https://en.wikipedia.org/wiki/Ordinal_definable_set})
\item Article : W. H. Woodin, "In search of Ultimate-L: the 19th Midrasha Mathematicae Lectures" (2017), Bull. Symbolic Logic 23
\end{refs}

\bigskip

\begin{problem}[{La frontière de la décidabilité : le problème de la fonction exponentielle de Tarski --- Recherche}]
\textbf{FND-22.} \textbf{Problème.} Tarski a démontré (1951) que la théorie du premier ordre du
corps des réels $ (\mathbb{R};\, +, \cdot, \le) $ est décidable, par élimination des
quantificateurs.
\begin{enumerate}
\item[(a)] En guise d'échauffement, effectuez une étape d'élimination à la main : trouvez une condition
sans quantificateur sur les paramètres $ a, b, c $ équivalente sur $\mathbb{R}$ à
$ \exists x\, (a x^2 + b x + c = 0) $, et démontrez l'équivalence.
\item[(b)] La question de recherche --- le \textit{problème de la fonction exponentielle
de Tarski} : la théorie du premier ordre de $ (\mathbb{R};\, +, \cdot, \le, e^x) $
est-elle décidable ? \textit{Statut : ouvert. Macintyre et Wilkie ont démontré en 1996 que la
décidabilité découle de la conjecture de Schanuel en théorie des nombres transcendants, elle-même
grande ouverte ; Wilkie a démontré que la théorie est modèle-complète et o-minimale, de sorte que
la géométrie est apprivoisée --- seule manque la procédure de décision.}
\end{enumerate}

\end{problem}

\noindent Le contraste avec FND-19 et FND-16 situe la frontière exactement : la
multiplication entière donne l'indécidabilité, la multiplication réelle est décidable, et
l'exponentiation réelle est le poste-frontière ouvert. Le problème relie la logique à une théorie
des nombres profonde : un unique énoncé conjectural sur des degrés de transcendance trancherait
une question portant sur des algorithmes. L'o-minimalité, née de ce cercle d'idées, a ensuite
alimenté le comptage de points de Pila--Wilkie et ses applications à la géométrie
diophantienne.

\begin{refs}
\item Contexte : Problème de la fonction exponentielle de Tarski --- Wikipédia (en anglais) (\url{https://en.wikipedia.org/wiki/Tarski%27s_exponential_function_problem})
\item Contexte : Conjecture de Schanuel --- Wikipédia (\url{https://fr.wikipedia.org/wiki/Conjecture_de_Schanuel})
\item Article : A. Macintyre and A. J. Wilkie, "On the decidability of the real exponential field" (1996), in \textit{Kreiseliana: About and Around Georg Kreisel}
\end{refs}

\bigskip

\begin{problem}[{p = t : un problème de soixante-dix ans résolu --- Recherche}]
\textbf{FND-23.} \textbf{Problème.} Pour des ensembles infinis $ A, B \subseteq \mathbb{N} $, on écrit
$ A \subseteq^{*} B $ si $ A \setminus B $ est fini. Le \textit{nombre de
pseudo-intersection} $ \mathfrak{p} $ est le plus petit cardinal d'une famille $ F $ de
parties infinies de $\mathbb{N}$ telle que toute sous-famille finie ait une intersection infinie, sans
qu'aucun $ A $ infini ne vérifie $ A \subseteq^{*} B $ pour tout $ B \in F $. Le \textit{nombre
de tour} $ \mathfrak{t} $ est la plus petite longueur d'une suite $ \subseteq^{*} $-décroissante
de parties infinies de $\mathbb{N}$ telle qu'aucun $ A $ infini ne soit presque contenu dans toutes.
\begin{enumerate}
\item[(a)] Démontrez, dans ZFC, que $ \aleph_1 \le \mathfrak{p} \le \mathfrak{t} \le 2^{\aleph_0} $.
\item[(b)] Savoir si $ \mathfrak{p} < \mathfrak{t} $ est cohérent est resté ouvert quelque soixante-dix
ans. \textit{Statut : résolu --- Malliaris et Shelah ont démontré $ \mathfrak{p} = \mathfrak{t} $
dans ZFC (annoncé en 2012--13, publié en 2016), par une voie inattendue passant par l'ordre de
Keisler, issu de la théorie des modèles.} Étudiez leur stratégie de démonstration et expliquez
pourquoi une question portant sur $\mathbb{R}$ a reçu une réponse par la théorie des modèles.
\end{enumerate}

\end{problem}

\noindent Les caractéristiques cardinales du continu mesurent les nombreuses tailles
distinctes que peuvent avoir les \og{} petites \fg{} parties de $\mathbb{R}$ quand HC est fausse ; presque toutes les
paires de caractéristiques avaient été séparées par forcing, et $ \mathfrak{p} $ contre
$ \mathfrak{t} $ était le survivant obstiné d'un travail remontant à Rothberger (1948). Le
théorème de Malliaris--Shelah --- récompensé par la médaille Hausdorff 2017 --- est l'exemple moderne
type d'un énoncé \og{} manifestement indépendant \fg{} se révélant être un théorème de ZFC : un
avertissement et un encouragement pour tout le reste de cette page.

\begin{refs}
\item Contexte : Caractéristique cardinale du continu --- Wikipédia (en anglais) (\url{https://en.wikipedia.org/wiki/Cardinal_characteristic_of_the_continuum})
\item Article : M. Malliaris and S. Shelah, "General topology meets model theory, on p and t" (2013), Proc. Nat. Acad. Sci. USA 110
\item Article : M. Malliaris and S. Shelah, "Cofinality spectrum theorems in model theory, set theory, and general topology" (2016), J. Amer. Math. Soc. 29
\end{refs}

\bigskip

\begin{problem}[{Le dixième problème de Hilbert sur les rationnels --- Recherche, short}]
\textbf{FND-33.} \textbf{Problème.} (Ouvert.) Sur les entiers, il n'existe aucun algorithme
décidant si une équation polynomiale à coefficients entiers a une solution (FND-31). Tranchez
l'analogue rationnel : existe-t-il un algorithme qui, étant donné un polynôme $ P $ à
coefficients rationnels, détermine si $ P(x_1, \ldots, x_n) = 0 $ a une solution dans $\mathbb{Q}$ ?
\end{problem}

\noindent La voie classique vers une réponse négative consiste à montrer que $\mathbb{Z}$ est
diophantien dans $\mathbb{Q}$, ce qui transférerait vers le bas l'indécidabilité de Matiiassevitch ; Jochen
Koenigsmann s'en est approché en 2016 en donnant une définition du premier ordre purement
universelle de $\mathbb{Z}$ à l'intérieur de $\mathbb{Q}$, mais une définition universelle n'est pas une définition
diophantienne. Le problème a beaucoup bougé sur le terrain voisin : en 2024--25, Peter Koymans et
Carlo Pagano, et indépendamment Levent Alpöge, Manjul Bhargava, Wei Ho et Ari Shnidman, ont établi
la réponse négative pour l'anneau des entiers de tout corps de nombres, en utilisant la stabilité
du rang des courbes elliptiques dans les extensions quadratiques. En 2026, le cas de $\mathbb{Q}$ lui-même
reste ouvert.

\begin{refs}
\item Contexte : Dixième problème de Hilbert --- Wikipédia (\url{https://fr.wikipedia.org/wiki/Dixi%C3%A8me_probl%C3%A8me_de_Hilbert})
\item Article : J. Koenigsmann, "Defining Z in Q" (2016), \textit{Annals of Mathematics} 183
\item Article : P. Koymans and C. Pagano, "Hilbert's tenth problem via additive combinatorics" (2024), arXiv:2412.01768 (\url{https://arxiv.org/abs/2412.01768})
\item Voir aussi : OPEN-63 (\url{open-problems.html#OPEN-63}) --- la formulation en théorie des nombres, sur la page des Problèmes ouverts.
\end{refs}

\bigskip

\begin{problem}[{Quelle taille le continu en $ \aleph_\omega $ peut-il atteindre ? --- Recherche, short}]
\textbf{FND-34.} \textbf{Problème.} (Ouvert.) Supposons $ \aleph_\omega $ cardinal limite fort,
c'est-à-dire $ 2^{\aleph_n} \aleph_{\omega_1} $ --- de façon
équivalente, si $ \mathrm{pcf}\{\aleph_n : n < \omega\} $ peut être non dénombrable.
\end{problem}

\noindent Easton a montré en 1970 que ZFC ne dit presque rien de la fonction continu aux
cardinaux réguliers : elle peut être n'importe quelle fonction monotone respectant la cofinalité.
Les cardinaux singuliers se sont révélés différents. Le théorème de Silver de 1974, puis la
théorie pcf de Shelah (1989), ont extrait de véritables contraintes ZFC de la seule structure des
cofinalités, et la borne $ \aleph_{\omega_4} $ est l'emblème de cette théorie --- un nombre
étrange et spécifique que personne ne croit optimal et que personne ne sait améliorer. En 2026, la
question est ouverte, et c'est le problème central de l'arithmétique cardinale.

\begin{refs}
\item Contexte : Théorie pcf --- Wikipédia (en anglais) (\url{https://en.wikipedia.org/wiki/PCF_theory})
\item Contexte : Hypothèse des cardinaux singuliers --- Wikipédia (en anglais) (\url{https://en.wikipedia.org/wiki/Singular_cardinals_hypothesis})
\item Lecture : Saharon Shelah, \textit{Cardinal Arithmetic} (Oxford, 1994)
\end{refs}

\bigskip

\begin{problem}[{Quelle est la force du théorème de Hindman ? --- Recherche}]
\textbf{FND-35.} \textbf{Problème.} (Partiellement résolu ; la question centrale est ouverte.) Le théorème de
Hindman dit que, pour tout coloriage de $\mathbb{N}$ avec un nombre fini de couleurs, il existe un ensemble
infini $ X \subseteq \mathbb{N} $ tel que toute somme d'un nombre fini d'éléments distincts de
$ X $ reçoive la même couleur. Les mathématiques à rebours mesurent un théorème par le plus
faible sous-système de l'arithmétique du second ordre qui le démontre : ici, les systèmes
pertinents sont $ \mathrm{RCA}_0 $ (la théorie de base), $ \mathrm{ACA}_0 $ (clôture par saut
de Turing) et $ \mathrm{ACA}_0^{+} $ (clôture par le $ \omega $-ième saut de Turing).
\begin{enumerate}
\item[(a)] Démontrez le théorème de Hindman par l'argument de Galvin--Glazer : munissez les ultrafiltres
de $\mathbb{N}$ (FND-28) d'une opération de semi-groupe, montrez qu'elle possède un idempotent, et lisez le
théorème sur cet idempotent.
\item[(b)] Travaillez la démonstration de Blass, Hirst et Simpson selon laquelle le théorème de Hindman
implique $ \mathrm{ACA}_0 $ sur $ \mathrm{RCA}_0 $, et qu'il est démontrable dans
$ \mathrm{ACA}_0^{+} $.
\item[(c)] Expliquez pourquoi la démonstration par ultrafiltres de (a) est bien plus forte qu'elle n'en a
l'air de l'extérieur, et ce qu'une démonstration menée à l'intérieur de $ \mathrm{ACA}_0 $
devrait éviter.
\item[(d)] Le problème ouvert : cerner la force exacte. Le théorème de Hindman est-il équivalent à
$ \mathrm{ACA}_0 $, à $ \mathrm{ACA}_0^{+} $, ou strictement entre les deux ? Une question
ouverte apparentée : la restriction aux sommes d'au plus deux éléments distincts est-elle déjà
équivalente au théorème complet ?
\end{enumerate}

\end{problem}

\noindent Neil Hindman a démontré le théorème des sommes finies en 1974 par un long
argument combinatoire ; la démonstration par ultrafiltres de Glazer, s'appuyant sur Galvin, l'a
réduite à une demi-page et en a fait l'exemple fondateur de l'algèbre du compactifié de
Stone--Čech. Blass, Hirst et Simpson l'ont coincé entre $ \mathrm{ACA}_0 $ et
$ \mathrm{ACA}_0^{+} $ en 1987, et il y est resté : cet écart est depuis près de quarante ans
l'un des problèmes ouverts les plus connus des mathématiques à rebours, résistant à la machinerie
de calculabilité bâtie autour du théorème de Ramsey pour les paires (FND-27). En 2026, il est
ouvert.

\begin{refs}
\item Contexte : Théorème de Hindman --- Wikipédia (en anglais) (\url{https://en.wikipedia.org/wiki/Hindman%27s_theorem})
\item Contexte : Mathématiques à rebours --- Wikipédia (\url{https://fr.wikipedia.org/wiki/Math%C3%A9matiques_%C3%A0_rebours})
\item Lecture : Stephen G. Simpson, \textit{Subsystems of Second Order Arithmetic}
\item Article : N. Hindman, "Finite sums from sequences within cells of a partition of N" (1974), \textit{Journal of Combinatorial Theory, Series A} 17
\end{refs}

\bigskip

\begin{problem}[{Les Cinq Grands, et pourquoi ils sont cinq --- Recherche}]
\textbf{FND-44.} \textbf{Problème.} Les mathématiques à rebours demandent, d'un théorème des mathématiques
ordinaires, quels axiomes d'existence d'ensembles sont nécessaires pour le démontrer --- et
répondent en démontrant le théorème équivalent, sur une théorie de base faible, à ces axiomes. Les
cinq sous-systèmes récurrents de l'arithmétique du second ordre sont
$ \mathrm{RCA}_0 \subsetneq \mathrm{WKL}_0 \subsetneq \mathrm{ACA}_0 \subsetneq \mathrm{ATR}_0
\subsetneq \Pi^1_1\text{-}\mathrm{CA}_0 $.
\begin{enumerate}
\item[(a)] Démontrez deux équivalences sur $ \mathrm{RCA}_0 $. D'abord : $ \mathrm{WKL}_0 $ est
équivalent au théorème de Heine--Borel pour $ [0,1] $, et à l'énoncé que toute fonction réelle
continue sur $ [0,1] $ est uniformément continue. Ensuite : $ \mathrm{ACA}_0 $ est équivalent
au théorème de Bolzano--Weierstrass, et à l'énoncé que l'image de toute injection
$ \mathbb{N} \to \mathbb{N} $ existe en tant qu'ensemble.
\item[(b)] Démontrez que $ \mathrm{RCA}_0 $ ne démontre pas $ \mathrm{WKL}_0 $ : montrez que le
$ \omega $-modèle dont les ensembles sont exactement les ensembles calculables satisfait
$ \mathrm{RCA}_0 $, et exhibez un arbre binaire calculable infini sans chemin infini
calculable.
\item[(c)] Énoncez soigneusement le \og{} phénomène des Cinq Grands \fg{} : pour l'essentiel, tout théorème des
mathématiques classiques analysé jusqu'ici est soit démontrable dans $ \mathrm{RCA}_0 $, soit
équivalent sur cette base à l'un des quatre autres systèmes. Passez ensuite en revue les
échappées connues, toutes situées sous $ \mathrm{ACA}_0 $ : le théorème de Ramsey pour les
paires $ \mathrm{RT}^2_2 $ (FND-27) est strictement plus faible que $ \mathrm{ACA}_0 $ par le
théorème de Seetapun, et est incomparable avec $ \mathrm{WKL}_0 $ --- aucun des deux ne démontre
l'autre --- de sorte qu'il n'est équivalent à aucun des cinq. Étudiez le \og{} zoo des mathématiques à
rebours \fg{}, l'ensemble de principes qui a poussé dans cet intervalle.
\item[(d)] Le problème de recherche : \textbf{existe-t-il un théorème expliquant le phénomène des Cinq
Grands ?} Donnez une conjecture précise --- une propriété des énoncés mathématiques
\og{} naturels \fg{} sous laquelle les cinq systèmes épuisent démontrablement les possibilités --- ou
expliquez pourquoi le verdict honnête est que la régularité est empirique. \textit{Statut : ouvert, et
le formuler précisément fait partie du problème.}
\end{enumerate}

\end{problem}

\noindent Harvey Friedman a lancé le programme au Congrès international de 1974 et Stephen
Simpson en a fait une discipline ; la surprise fondatrice fut le peu d'axiomes dont les
mathématiques ordinaires ont réellement besoin, et la surprise plus profonde fut que tant de
théorèmes sans rapport atterrissent exactement sur les mêmes cinq barreaux. Personne ne sait
pourquoi. La meilleure preuve que le motif n'est pas un théorème en attente d'être trouvé est la
région sous $ \mathrm{ACA}_0 $ : Specker a montré en 1971 que le théorème de Ramsey pour les
paires admet des coloriages calculables sans ensemble homogène calculable, Seetapun et Slaman ont
démontré en 1995 qu'il reste néanmoins en deçà de $ \mathrm{ACA}_0 $, et Liu Jiayi a démontré en
2012 qu'il n'implique même pas le lemme de König faible. Ce qui y a poussé depuis est un véritable
zoo --- des dizaines de principes, deux à deux incomparables, dont aucun n'est l'un des cinq.
Comparez avec FND-35, où la force exacte d'un seul théorème est ouverte depuis près de quarante
ans. En 2026, le phénomène des Cinq Grands n'a ni démonstration ni explication admise.

\begin{refs}
\item Contexte : Mathématiques à rebours --- Wikipédia (\url{https://fr.wikipedia.org/wiki/Math%C3%A9matiques_%C3%A0_rebours})
\item Contexte : Théorème de Ramsey --- Wikipédia (\url{https://fr.wikipedia.org/wiki/Th%C3%A9or%C3%A8me_de_Ramsey}) (voir sa section sur les mathématiques à rebours)
\item Lecture : Stephen G. Simpson, \textit{Subsystems of Second Order Arithmetic}
\item Article : A. Montalbán, "Open questions in reverse mathematics", \textit{Bulletin of Symbolic Logic} 17 (2011), 431--454
\end{refs}

\bigskip

\begin{problem}[{Quelle est la force de l'axiome du forcing propre ? --- Recherche, short}]
\textbf{FND-45.} \textbf{Problème.} (\textbf{Ouvert.}) L'axiome du forcing propre PFA étend
l'axiome de Martin (FND-42) des ordres partiels ccc à tous les ordres \textit{propres} ; il implique
$ 2^{\aleph_0} = \aleph_2 $, et il est vrai dans une extension générique dès qu'il existe un
cardinal supercompact, tandis que la meilleure minoration connue de sa force de cohérence se situe
très en dessous, un peu en deçà d'un cardinal de Woodin limite de cardinaux de Woodin. Déterminez
la force de cohérence exacte de PFA --- un cardinal supercompact est-il vraiment nécessaire, ou la
force véritable est-elle bien moindre ?
\end{problem}

\noindent PFA est issu des travaux sur le forcing propre itéré du début des années 1980,
et c'est le plus fort des axiomes de forcing d'usage courant : contrairement à HC ou à sa négation,
il tranche un éventail énorme de questions combinatoires, et il tranche le continu. L'écart entre
sa majoration supercompacte et la minoration connue est probablement le plus large de ce genre en
théorie des ensembles et l'un des problèmes ouverts les plus cités du domaine. Il se trouve aussi
au centre d'un désaccord vivant : les axiomes de forcing donnent $ 2^{\aleph_0} = \aleph_2 $, le
programme Ultimate-L de Woodin (FND-21) pointe dans l'autre direction, et la quantité de force de
grand cardinal que coûte vraiment PFA fait partie de ce dont on débat. Voyez aussi FND-20 sur la
linéarité inexpliquée de la hiérarchie des grands cardinaux.

\begin{refs}
\item Contexte : Axiome du forcing propre --- Wikipédia (en anglais) (\url{https://en.wikipedia.org/wiki/Proper_forcing_axiom})
\item Contexte : Grand cardinal --- Wikipédia (\url{https://fr.wikipedia.org/wiki/Grand_cardinal})
\item Article : J. T. Moore, "The proper forcing axiom", \textit{Proceedings of the International Congress of Mathematicians} (Hyderabad, 2010)
\end{refs}

\bigskip


\section*{Pour aller plus loin}


\vfill
\noindent\emph{Hors de la Caverne} --- des probl\`emes, pas de r\'eponses.
\end{document}
